%%%%%%%%%%%%%%%%%%%%%%%%%%%%%%%%%%%%%%%%%%%%%%%%%%%%%%%%%%%%%%%%%%%%%%%%%%%%%%%
%  SÍLABO DE ECONOMÍA PÚBLICA — Plan de estudio integrado
%
%  Teoría de la economía pública:
%    GR  — Gruber, Public Finance and Public Policy (Worth, 3.ª ed., 2011)
%    SR  — Stiglitz y Rosengard, Economics of the Public Sector (Norton,
%          4.ª ed., 2015)
%    HM  — Hindriks y Myles, Intermediate Public Economics (borrador de
%          junio de 2004; MIT Press, 2006)
%    AS  — Atkinson y Stiglitz, Lectures on Public Economics (Princeton,
%          edición actualizada, 2015)
%    SA  — Salanié, The Economics of Taxation (MIT Press, 2003)
%  Elección pública, evaluación y finanzas públicas:
%    MU  — Mueller, Public Choice III (Cambridge University Press, 2003)
%    DR  — Drazen, Political Economy in Macroeconomics (Princeton, 2000)
%    BO  — Boardman, Greenberg, Vining y Weimer, Cost-Benefit Analysis
%          (Cambridge University Press, 5.ª ed., 2018)
%    BG  — Barrantes, Manrique y Glave, Economía pública (Material de
%          Enseñanza 9, PUCP, 2.ª ed., 2023)
%  Textos complementarios con horas:
%    PT  — Persson y Tabellini, Political Economics (MIT Press, 2000)
%    HE  — Heer, Public Economics: The Macroeconomic Perspective (Springer, 2019)
%    KC  — Kocherlakota, The New Dynamic Public Finance (Princeton, 2010)
%    MA  — Mendoza y Anastacio, La historia fiscal del Perú: 1980-2020
%          (Fondo Editorial PUCP, 2021)
%    RB  — Corbacho, Fretes Cibils y Lora (eds.), Recaudar no basta (BID, 2013)
%    UR  — Urrunaga, Hiraoka y Risso, Fundamentos de economía pública
%          (Universidad del Pacífico, 2014)
%    NC  — Arias, notas de clase (diapositivas) de 1ECO36, Finanzas Públicas
%  Programas integrados: 1ECO18 (Economía Pública), 1ECO36 (Finanzas
%    Públicas), 1ECO81 (Economía Pública, maestría) y 1ECO40 (Pensiones y
%    Seguridad Social) de la PUCP; MIT OCW 14.41 (otoño de 2024) y ECO1028
%    (Public Choice, Dublin City University, 2026) como referencias externas.
%
%  Compilar con:  pdflatex EconomiaPublica.tex   (dos veces, por el índice)
%%%%%%%%%%%%%%%%%%%%%%%%%%%%%%%%%%%%%%%%%%%%%%%%%%%%%%%%%%%%%%%%%%%%%%%%%%%%%%%

\documentclass[11pt,a4paper]{article}

\usepackage[T1]{fontenc}
\usepackage[utf8]{inputenc}
\usepackage[spanish,es-noquoting,es-tabla]{babel}
\usepackage{lmodern}
\usepackage[scaled=0.92]{helvet}        % títulos en sans serif, al estilo McGraw-Hill
\usepackage[margin=2.4cm,headheight=15pt]{geometry}
\usepackage[table]{xcolor}
\usepackage{enumitem}
\usepackage{titlesec}
\usepackage{titletoc}
\usepackage{tabularx}
\usepackage{booktabs}
\usepackage{longtable}
\usepackage{amsmath}
\usepackage{amssymb}
\usepackage{fancyhdr}
\usepackage{tikz}
\usepackage[skins]{tcolorbox}
\usepackage{needspace}

%------------------------------------------------------------------ colores ---
% Gama de naranjas
\definecolor{tonoOscuro}{RGB}{82,36,4}              % oscuro
\definecolor{tonoProfundo}{RGB}{121,52,6}           % profundo
\definecolor{tonoPrincipal}{RGB}{160,69,8}          % principal
\definecolor{tonoA}{RGB}{179,117,25}                % variante cálida
\definecolor{tonoB}{RGB}{179,53,15}                 % variante fría
\definecolor{tonoC}{RGB}{193,33,21}                 % variante viva
\definecolor{tonoClaro}{RGB}{221,146,95}            % claro
\definecolor{tonoAcero}{RGB}{131,90,63}             % acero
\definecolor{tonoPizarra}{RGB}{134,105,86}          % pizarra
\definecolor{tonoGrafito}{RGB}{73,53,39}            % grafito (texto secundario)
\definecolor{tonoHielo}{RGB}{247,223,207}           % hielo
\definecolor{tonoPalido}{RGB}{251,240,233}          % pálido
\definecolor{tonoBruma}{RGB}{254,249,246}           % bruma

\colorlet{azulosc}{tonoOscuro}
\colorlet{grisclaro}{tonoPalido}
% Un tono distinto de la gama para cada libro
\definecolor{colorGR}{RGB}{124,17,14}
\definecolor{colorSR}{RGB}{170,32,19}
\definecolor{colorHM}{RGB}{147,34,16}
\definecolor{colorAS}{RGB}{124,35,14}
\definecolor{colorSA}{RGB}{170,56,19}
\definecolor{colorMU}{RGB}{147,55,16}
\definecolor{colorDR}{RGB}{124,52,14}
\definecolor{colorBO}{RGB}{170,79,19}
\definecolor{colorBG}{RGB}{147,75,16}
\definecolor{colorPT}{RGB}{124,69,14}
\definecolor{colorHE}{RGB}{170,103,19}
\definecolor{colorKC}{RGB}{147,96,16}
\definecolor{colorMA}{RGB}{124,87,14}
\definecolor{colorRB}{RGB}{170,127,19}
\definecolor{colorUR}{RGB}{147,116,16}
\colorlet{colorPUCP}{tonoPizarra}

\usepackage[colorlinks=true,linkcolor=tonoOscuro,urlcolor=tonoPrincipal]{hyperref}

%------------------------------------------------------------------ estilos ---
\titleformat{\section}[block]
  {\sffamily\LARGE\bfseries\color{tonoOscuro}\raggedright}{\thesection.}{0.6em}
  {\textcolor{tonoPrincipal}{\rule[-0.1ex]{0.55em}{0.85em}}\hspace{0.45em}}
  [\vspace{-0.4ex}{\color{tonoPrincipal}\titlerule[1.6pt]}\vspace{1.2pt}{\color{tonoClaro}\titlerule[0.6pt]}]
\titlespacing*{\section}{0pt}{2.6ex plus 1ex minus .2ex}{1.6ex plus .2ex}

\titleformat{\subsection}
  {\sffamily\large\bfseries\color{tonoPrincipal}}{\thesubsection}{0.6em}
  {\textcolor{tonoClaro}{\rule[0.1ex]{0.45em}{0.45em}}\hspace{0.45em}}
\titlespacing*{\subsection}{0pt}{2.2ex plus .8ex minus .2ex}{1ex plus .2ex}

% Índice
\titlecontents{section}[0pt]
  {\addvspace{3pt}\sffamily\bfseries\color{tonoOscuro}}
  {}{\textcolor{tonoB}{\rule[0.15ex]{0.4em}{0.4em}}\hspace{0.6em}}
  {{\color{tonoClaro}\titlerule*[0.7em]{.}}\textcolor{tonoPrincipal}{\contentspage}}

% Cabecera y pie
\pagestyle{fancy}
\fancyhf{}
\lhead{\small\sffamily\bfseries\color{tonoOscuro}Sílabo de Economía Pública}
\rhead{\footnotesize\sffamily\color{tonoPrincipal}Gruber $\cdot$ Stiglitz--Rosengard $\cdot$ Atkinson--Stiglitz $\cdot$ Mueller}
\cfoot{{\setlength{\fboxsep}{3pt}\colorbox{tonoPrincipal}{\makebox[1.8em]{\small\sffamily\bfseries\color{white}\thepage}}}}
\renewcommand{\headrule}{%
  {\color{tonoPrincipal}\hrule height 1.1pt width\headwidth}\vskip 1.2pt
  {\color{tonoClaro}\hrule height 0.4pt width\headwidth}\vskip -2.7pt}

% Listas
\setlist[itemize,1]{label=\textcolor{tonoPrincipal}{\scriptsize$\blacksquare$}}
\setlist[enumerate,1]{label=\textcolor{tonoPrincipal}{\sffamily\bfseries\arabic*.}}

% Tablas: cabecera oscura y filas alternas en tono pálido
\arrayrulecolor{tonoPrincipal}
\setlength{\aboverulesep}{0pt}
\setlength{\belowrulesep}{0pt}
\renewcommand{\arraystretch}{1.3}
\newcommand{\cab}[1]{\textcolor{white}{\sffamily\bfseries #1}}
\newcommand{\filasalternas}{\rowcolors{2}{white}{tonoPalido}}

% Encabezados de párrafo en color
\newcommand{\destacado}[1]{{\sffamily\bfseries\color{tonoPrincipal}#1}}

% Franja decorativa con la gama
\newcommand{\franjas}[1][3pt]{%
  \noindent\begin{tikzpicture}
    \foreach \c [count=\i from 0] in {tonoOscuro,tonoProfundo,tonoPrincipal,tonoA,tonoB,tonoC,tonoClaro,tonoHielo}
      \fill[\c] (\i*\linewidth/8,0) rectangle ++(\linewidth/8,#1);
  \end{tikzpicture}\par}

%---------------------------------------------------------------- comandos ----
% Etiqueta de libro: recuadro de color con la sigla
\newcommand{\etiqueta}[2]{%
  {\setlength{\fboxsep}{1.4pt}\colorbox{#1}{\sffamily\bfseries\color{white}\strut#2}}}
% Referencias abreviadas a cada libro
\newcommand{\GR}[1]{\etiqueta{colorGR}{GR}\nobreak\hspace{0.25em}\textcolor{colorGR}{#1}}
\newcommand{\SR}[1]{\etiqueta{colorSR}{SR}\nobreak\hspace{0.25em}\textcolor{colorSR}{#1}}
\newcommand{\HM}[1]{\etiqueta{colorHM}{HM}\nobreak\hspace{0.25em}\textcolor{colorHM}{#1}}
\newcommand{\AS}[1]{\etiqueta{colorAS}{AS}\nobreak\hspace{0.25em}\textcolor{colorAS}{#1}}
\newcommand{\SA}[1]{\etiqueta{colorSA}{SA}\nobreak\hspace{0.25em}\textcolor{colorSA}{#1}}
\newcommand{\MU}[1]{\etiqueta{colorMU}{MU}\nobreak\hspace{0.25em}\textcolor{colorMU}{#1}}
\newcommand{\DR}[1]{\etiqueta{colorDR}{DR}\nobreak\hspace{0.25em}\textcolor{colorDR}{#1}}
\newcommand{\BO}[1]{\etiqueta{colorBO}{BO}\nobreak\hspace{0.25em}\textcolor{colorBO}{#1}}
\newcommand{\BG}[1]{\etiqueta{colorBG}{BG}\nobreak\hspace{0.25em}\textcolor{colorBG}{#1}}
\newcommand{\PT}[1]{\etiqueta{colorPT}{PT}\nobreak\hspace{0.25em}\textcolor{colorPT}{#1}}
\newcommand{\HE}[1]{\etiqueta{colorHE}{HE}\nobreak\hspace{0.25em}\textcolor{colorHE}{#1}}
\newcommand{\KC}[1]{\etiqueta{colorKC}{KC}\nobreak\hspace{0.25em}\textcolor{colorKC}{#1}}
\newcommand{\MA}[1]{\etiqueta{colorMA}{MA}\nobreak\hspace{0.25em}\textcolor{colorMA}{#1}}
\newcommand{\RB}[1]{\etiqueta{colorRB}{RB}\nobreak\hspace{0.25em}\textcolor{colorRB}{#1}}
\newcommand{\UR}[1]{\etiqueta{colorUR}{UR}\nobreak\hspace{0.25em}\textcolor{colorUR}{#1}}
\newcommand{\NC}[1]{\textcolor{colorPUCP}{\textbf{NC#1}}}

% Un tema: número en recuadro de color + título
\newcommand{\tema}[1]{\temaaux#1\relax}
\def\temaaux#1 #2\relax{%
  \par\addvspace{1.1em}\noindent
  {\sffamily\bfseries\large\raggedright
   \hangindent=2.9em\hangafter=1
   \makebox[2.9em][l]{{\setlength{\fboxsep}{2.5pt}\colorbox{tonoPrincipal}{\color{white}#1}}}%
   {\color{tonoOscuro}#2}\par}%
  \nobreak\vspace{0.25em}\noindent{\color{tonoClaro}\rule{\linewidth}{0.6pt}}\par\nobreak\vspace{0.35em}}

% Línea de lecturas con marcador y etiqueta de color
\newcommand{\linealectura}[3]{%
  \par\nobreak\noindent\hspace{1.2em}%
  \begin{minipage}{0.92\textwidth}\small
    \textcolor{#1}{\rule[0.1ex]{0.45em}{0.45em}}\hspace{0.45em}%
    {\sffamily\bfseries\color{#1}#2}\enspace#3%
  \end{minipage}\par\vspace{0.3em}}
\newcommand{\lecturas}[1]{\linealectura{tonoPrincipal}{Lecturas:}{#1}}
\newcommand{\apoyo}[1]{\linealectura{tonoA}{Profundización:}{#1}}
\newcommand{\juegos}[1]{\linealectura{tonoB}{Teoría de juegos:}{#1}}
\newcommand{\supl}[1]{\linealectura{tonoC}{Suplemento:}{#1}}
\newtcolorbox{cajanota}{enhanced, sharp corners, boxrule=0pt, frame hidden,
  colback=tonoPalido, borderline west={2.5pt}{0pt}{tonoClaro},
  left=8pt, right=6pt, top=4pt, bottom=4pt, before skip=0pt, after skip=0pt,
  fontupper=\small\color{tonoGrafito}}
\newcommand{\nota}[1]{%
  \par\nobreak\noindent\hspace{1.2em}%
  \begin{minipage}{0.92\textwidth}\begin{cajanota}#1\end{cajanota}\end{minipage}\par\vspace{0.3em}}
\newcommand{\consultalib}[1]{\linealectura{tonoPizarra}{Consulta:}{{\color{tonoGrafito}#1}}}
% Artículos y textos tomados de los programas PUCP
\newcommand{\articulos}[1]{\linealectura{colorPUCP}{Artículos (PUCP):}{#1}}
\newcommand{\consulta}[1]{\linealectura{colorPUCP}{Consulta (PUCP):}{{\color{tonoGrafito}#1}}}

% Encabezado de bloque: banda oscura, franja de la gama y recuadro de presentación
\newif\ifrecienparte
\newcommand{\bloque}[2]{%
  \ifrecienparte\global\recienpartefalse\else\par\needspace{0.4\textheight}\fi
  \phantomsection\addcontentsline{toc}{section}{#1}%
  \begin{tcolorbox}[enhanced, sharp corners, boxrule=0pt, frame hidden,
      colback=tonoOscuro, borderline west={7pt}{0pt}{tonoClaro},
      left=14pt, right=10pt, top=10pt, bottom=10pt, before skip=0pt, after skip=0pt]
    {\raggedright\sffamily\Large\bfseries\color{white}#1\par}
  \end{tcolorbox}%
  \nointerlineskip\franjas[4pt]
  \vspace{0.8em}
  \begin{tcolorbox}[enhanced, sharp corners, boxrule=0pt, frame hidden,
      colback=tonoPalido, borderline west={3pt}{0pt}{tonoPrincipal},
      borderline south={0.6pt}{0pt}{tonoClaro},
      left=10pt, right=8pt, top=6pt, bottom=6pt, before skip=0pt, after skip=0pt,
      fontupper=\small\color{tonoGrafito}]
    #2
  \end{tcolorbox}%
  \par\medskip}

% Encabezado de parte: banda a todo el ancho en el tono principal
\newcommand{\parte}[1]{%
  \clearpage
  \phantomsection\addcontentsline{toc}{section}{\textbf{#1}}%
  \begin{tcolorbox}[enhanced, sharp corners, boxrule=0pt, frame hidden,
      colback=tonoPrincipal, halign=center, top=16pt, bottom=16pt,
      borderline north={3pt}{0pt}{tonoClaro}, borderline south={3pt}{0pt}{tonoClaro},
      before skip=0pt, after skip=0pt]
    {\sffamily\LARGE\bfseries\color{white}#1}
  \end{tcolorbox}\par\bigskip
  \global\recienpartetrue}

%================================================================== documento ==
\begin{document}

\begin{titlepage}
\begin{tikzpicture}[remember picture,overlay]
  % banda superior
  \fill[tonoOscuro] (current page.north west) rectangle ([yshift=-5.4cm]current page.north east);
  \fill[tonoClaro] (current page.north west) rectangle ([xshift=1.1cm,yshift=-5.4cm]current page.north west);
  % franja de la gama bajo la banda
  \foreach \c [count=\i from 0] in {tonoOscuro,tonoProfundo,tonoPrincipal,tonoA,tonoB,tonoC,tonoClaro,tonoHielo}
    \fill[\c] ([xshift=\i*\paperwidth/8,yshift=-5.4cm]current page.north west)
      rectangle ++(\paperwidth/8,-0.35cm);
  % franja inferior
  \foreach \c [count=\i from 0] in {tonoHielo,tonoClaro,tonoC,tonoB,tonoA,tonoPrincipal,tonoProfundo,tonoOscuro}
    \fill[\c] ([xshift=\i*\paperwidth/8]current page.south west)
      rectangle ++(\paperwidth/8,0.6cm);
\end{tikzpicture}%
\centering
\vspace*{0.5cm}
{\sffamily\Huge\bfseries\color{white} Sílabo de Economía Pública\par}
\vspace{0.3cm}
{\sffamily\Large\color{tonoHielo} Plan de estudio integrado\par}
\vspace{1.9cm}
{\sffamily\large\bfseries\color{tonoPrincipal}\textcolor{tonoClaro}{\rule[0.5ex]{1.5cm}{0.6pt}}\quad Teoría de la economía pública\quad\textcolor{tonoClaro}{\rule[0.5ex]{1.5cm}{0.6pt}}\par}
\vspace{0.25cm}
{\large\etiqueta{colorGR}{GR} \quad {\color{tonoOscuro}Jonathan Gruber}\par}
{\normalsize \emph{\color{colorGR}Public Finance and Public Policy}\par}
{\small\color{tonoPizarra} Worth Publishers, 3.ª edición, 2011 --- 25 capítulos\par}
\vspace{0.15cm}
{\large\etiqueta{colorSR}{SR} \quad {\color{tonoOscuro}Joseph E. Stiglitz y Jay K. Rosengard}\par}
{\normalsize \emph{\color{colorSR}Economics of the Public Sector}\par}
{\small\color{tonoPizarra} W.\,W. Norton, 4.ª edición, 2015 --- 28 capítulos\par}
\vspace{0.15cm}
{\large\etiqueta{colorHM}{HM} \quad {\color{tonoOscuro}Jean Hindriks y Gareth D. Myles}\par}
{\normalsize \emph{\color{colorHM}Intermediate Public Economics}\par}
{\small\color{tonoPizarra} Borrador de junio de 2004 (The MIT Press, 2006) --- 23 capítulos\par}
\vspace{0.15cm}
{\large\etiqueta{colorAS}{AS} \quad {\color{tonoOscuro}Anthony B. Atkinson y Joseph E. Stiglitz}\par}
{\normalsize \emph{\color{colorAS}Lectures on Public Economics}\par}
{\small\color{tonoPizarra} Princeton University Press, edición actualizada, 2015 --- 18 lecciones\par}
\vspace{0.15cm}
{\large\etiqueta{colorSA}{SA} \quad {\color{tonoOscuro}Bernard Salanié}\par}
{\normalsize \emph{\color{colorSA}The Economics of Taxation}\par}
{\small\color{tonoPizarra} The MIT Press, 2003 --- 10 capítulos\par}
\vspace{0.15cm}
{\large\etiqueta{colorBG}{BG} \quad {\color{tonoOscuro}Roxana Barrantes, Silvana Manrique y Carla Glave}\par}
{\normalsize \emph{\color{colorBG}Economía pública}\par}
{\small\color{tonoPizarra} Material de Enseñanza 9, Departamento de Economía de la PUCP, 2.ª edición, 2023 --- 12 capítulos\par}
\vspace{0.4cm}
{\sffamily\large\bfseries\color{tonoPrincipal}\textcolor{tonoClaro}{\rule[0.5ex]{1.5cm}{0.6pt}}\quad Elección pública, economía política y evaluación\quad\textcolor{tonoClaro}{\rule[0.5ex]{1.5cm}{0.6pt}}\par}
\vspace{0.25cm}
{\large\etiqueta{colorMU}{MU} \quad {\color{tonoOscuro}Dennis C. Mueller}\par}
{\normalsize \emph{\color{colorMU}Public Choice III}\par}
{\small\color{tonoPizarra} Cambridge University Press, 2003 --- 29 capítulos\par}
\vspace{0.15cm}
{\large\etiqueta{colorDR}{DR} \quad {\color{tonoOscuro}Allan Drazen}\par}
{\normalsize \emph{\color{colorDR}Political Economy in Macroeconomics}\par}
{\small\color{tonoPizarra} Princeton University Press, 2000 --- 14 capítulos\par}
\vspace{0.15cm}
{\large\etiqueta{colorBO}{BO} \quad {\color{tonoOscuro}Anthony E. Boardman, David H. Greenberg, Aidan R. Vining y David L. Weimer}\par}
{\normalsize \emph{\color{colorBO}Cost--Benefit Analysis: Concepts and Practice}\par}
{\small\color{tonoPizarra} Cambridge University Press, 5.ª edición, 2018 --- 20 capítulos\par}
\vspace{0.4cm}
{\normalsize\color{colorPUCP} Con temas y lecturas de cuatro programas de la PUCP y de MIT 14.41,\par}
{\normalsize\color{colorPUCP} y con informes del MEF, la OCDE, el Banco Mundial, el BID y la CEPAL\par}
\vfill
{\setlength{\fboxsep}{6pt}\colorbox{tonoPrincipal}{\normalsize\sffamily\bfseries\color{white} Catorce bloques temáticos --- 7 meses y medio}\par}
\vspace{0.1cm}
\end{titlepage}

\tableofcontents
\newpage

%------------------------------------------------------------------------------
\section*{Cómo usar este sílabo}
\addcontentsline{toc}{section}{Cómo usar este sílabo}

El plan está organizado en \textbf{catorce bloques temáticos}, agrupados en
cuatro partes. La \textbf{Parte I} (Bloques 1 a 4) da los fundamentos: el
sector público y sus herramientas, las externalidades, los bienes públicos y
el federalismo fiscal, y la elección pública. La \textbf{Parte II} (Bloques 5
a 7) estudia el gasto: su análisis y evaluación, el seguro social y las
pensiones, la salud, la educación y la redistribución. La \textbf{Parte III}
(Bloques 8 a 10) es la tributación: incidencia y eficiencia, tributación
óptima y tributación aplicada. La \textbf{Parte IV} (Bloques 11 a 14) parte de
la economía política de la política fiscal y lleva todo lo anterior a las
finanzas públicas en la práctica, con el Perú como caso: cuentas fiscales, deuda y reglas, presupuesto y calidad del gasto,
descentralización, sistema tributario y reforma fiscal. Conviene cerrar cada
bloque antes de abrir el siguiente.

\medskip
\noindent\destacado{Por qué un curso aparte.} La economía pública aparece en los
demás sílabos solo en partes. El de microeconomía dedica su Bloque 11 a las
externalidades, los bienes públicos y el bienestar social, y su Bloque 13 al
diseño de mecanismos. El de macroeconomía trata la restricción
presupuestaria del gobierno, la equivalencia ricardiana, la sostenibilidad de
la deuda, el suavizamiento de impuestos y la economía política de los
déficits (temas 2.4 a 2.8), la imposición del capital en el modelo de Ramsey
(4.8), la seguridad social (7.3 y 7.5) y la teoría fiscal del nivel de precios
(8.7). Ninguno estudia la tributación (incidencia, exceso de gravamen,
tributación óptima y aplicada), el análisis y la evaluación del gasto, el
seguro social, la elección pública ni las instituciones de las finanzas
públicas peruanas, que son el núcleo de este plan.

\medskip
\noindent\destacado{Relación con los demás sílabos.} El plan supone la teoría
del consumidor y la dualidad (Bloque 2 del sílabo de microeconomía), el
equilibrio general (Bloque 9) y la información asimétrica (Bloque 12), y la
optimización estática del sílabo de economía matemática (Bloque 7). La
Parte III usa a fondo la dualidad y el principio de revelación (tema 13.3 del
sílabo de microeconomía). Los temas de evaluación (1.3 y 5.6) remiten al
sílabo de microeconometría para los métodos; los de economía política, deuda y
política fiscal (Bloques 11 y 12), a los Bloques 2, 9 y 13 del sílabo de
macroeconomía. Donde hay
superposición, lo ya leído se cita aquí como consulta.

\medskip
\noindent\destacado{Criterio general de reparto} (en el orden en que se lee cada tema):
\begin{itemize}[leftmargin=1.4em,itemsep=2pt]
  \item \textcolor{colorGR}{\textbf{Gruber}} abre la mayoría de los temas con
        la intuición, los gráficos y la evidencia empírica. Es el texto de
        MIT 14.41. Se lee \emph{primero}.
  \item \textcolor{colorSR}{\textbf{Stiglitz y Rosengard}} es el texto base de
        1ECO18: aporta la mirada institucional, el marco del análisis del
        gasto y el tratamiento de cada impuesto.
  \item \textcolor{colorBG}{\textbf{Barrantes, Manrique y Glave}} son las notas
        de 1ECO18: cortas, en castellano y con el caso peruano (sector
        público, análisis costo-beneficio de normas, análisis de impacto
        regulatorio, tributación y presupuesto).
  \item \textcolor{colorHM}{\textbf{Hindriks y Myles}} dan la versión formal de
        nivel intermedio de casi todos los temas teóricos y sirven de puente
        hacia los textos de posgrado.
  \item \textcolor{colorAS}{\textbf{Atkinson y Stiglitz}} y
        \textcolor{colorSA}{\textbf{Salanié}} son los textos de posgrado de la
        tributación (incidencia, exceso de gravamen y tributación óptima) y el
        núcleo del Módulo 3 de 1ECO81.
  \item \textcolor{colorMU}{\textbf{Mueller}} es el tratado de elección
        pública: reglas de votación, competencia electoral, búsqueda de
        rentas, burocracia, federalismo, tamaño del Estado y ciclos políticos
        (Bloque 4 y temas 3.4, 11.2 y 11.6).
  \item \textcolor{colorDR}{\textbf{Drazen}} es el texto de la economía
        política de la política fiscal: inconsistencia temporal y compromiso,
        ciclos políticos, redistribución, demoras de las reformas y sesgo hacia
        el déficit (Bloque 11 y tema 4.4).
  \item \textcolor{colorBO}{\textbf{Boardman et al.}} es el manual de análisis
        costo-beneficio (Bloque 5).
  \item \textcolor{colorPT}{\textbf{Persson y Tabellini}},
        \textcolor{colorHE}{\textbf{Heer}} y
        \textcolor{colorKC}{\textbf{Kocherlakota}} entran en dosis puntuales:
        la economía política formal, la macroeconomía de las finanzas
        públicas y la nueva finanza pública dinámica.
  \item Las \textcolor{colorPUCP}{\textbf{notas de clase de 1ECO36}}
        (Arias), \textcolor{colorRB}{\textbf{BID (2013)}},
        \textcolor{colorMA}{\textbf{Mendoza y Anastacio}} y los informes del
        MEF, la OCDE, el Banco Mundial y la CEPAL cubren las finanzas
        públicas peruanas y latinoamericanas (Parte IV y Bloque 10).
\end{itemize}

\medskip
\noindent\destacado{Qué entra en las horas.} Las líneas \emph{Lecturas},
\emph{Profundización} y \emph{Artículos (PUCP)} entran en el cómputo de horas,
incluidas las diapositivas de 1ECO36 y los informes de la carpeta. Las
líneas \emph{Consulta} y \emph{Consulta (PUCP)} remiten a libros, documentos
y referencias de los programas que no entran (Anexo E). Los informes extensos
se leen por las secciones indicadas; sus anexos estadísticos se consultan.

\medskip
\noindent\destacado{Sobre los ejercicios.} MIT 14.41 publica sus conjuntos de
problemas y sus exámenes con soluciones (Anexo D.1). Gruber cierra cada
capítulo con preguntas y problemas, y Hindriks y Myles y Persson y Tabellini
con ejercicios formales. Se recomienda reservar dos de las diez horas
semanales para ejercicios en las Partes I a III y, en la Parte IV, para
trabajar con las cifras del MEF (Consulta Amigable, cuadros del MMM) y del
BCRP.

\medskip
\noindent\destacado{Material excluido.} Quedan fuera los métodos numéricos (las
secciones cuantitativas de Heer y el capítulo 6 de Kocherlakota), la teoría de
la regulación y la organización industrial (1ECO82), las instituciones
políticas comparadas, la teoría monetaria y los modelos de proyección del
BCRP que hay en la carpeta, y la contabilidad nacional. El detalle está en el
Anexo B.

\medskip
\noindent\destacado{Notación.} \GR{n.m}, \HM{n.m}, \SA{n.m}, \MU{n.m}, \DR{n.m},
\PT{n.m}, \HE{n.m} y \MA{n.m} remiten a la sección $m$ del capítulo $n$ de cada libro, y
\BG{n.m}, a las notas de 1ECO18. Atkinson y Stiglitz numeran sus secciones
«12--2»: aquí se escriben \AS{12.2}. \SR{n}, \BO{n}, \KC{n} y \RB{n} remiten
a capítulos, porque sus índices no numeran las secciones. \UR{II.2.3} es la
sección 2.3 de la parte II de Urrunaga et al. \NC{3} es la nota de clase 3
(diapositivas) de 1ECO36, y «diap.» indica las diapositivas. Gruber se cita
por la 3.ª edición de la carpeta (2011); MIT 14.41 usa la 7.ª (2022), así que
con otra edición conviene ubicar los temas por título. El borrador de
Hindriks y Myles numera sus capítulos de otro modo que la edición de 2006:
la tributación de bienes, que aquí es \HM{15}, es el cap.~14 del libro que
cita 1ECO81.

%------------------------------------------------------------------------------
\newpage
\section*{Cronograma por semanas}
\addcontentsline{toc}{section}{Cronograma por semanas}

Calendario del plan: 330 horas distribuidas en \textbf{33 semanas de
contenido}, que equivalen a \textbf{unos 7 meses y medio de calendario}, a un
ritmo de diez horas semanales. El núcleo de 30 semanas (7 meses) se amplió
en tres para incorporar el Bloque 11, basado en Drazen.

\medskip
\noindent Todas las semanas son de materia nueva y no hay semanas de colchón.
La columna de horas es la que gobierna el reparto. Si hay semanas perdidas por
feriados, viajes o imprevistos, el ajuste correcto es alargar el calendario,
no comprimir los bloques.

\medskip
{\filasalternas
\noindent\begin{tabularx}{\textwidth}{lXcc}
\toprule
\rowcolor{tonoOscuro}\cab{Semanas} & \cab{Contenido} & \cab{Sem.} & \cab{Horas} \\
\midrule
\rowcolor{tonoA}\multicolumn{4}{l}{\cab{\textit{Parte I --- Fundamentos}}} \\
1--2   & Bloque 1 --- El sector público y las herramientas de la economía pública & 2 & 20 \\
3      & Bloque 2 --- Externalidades & 1 & 10 \\
4--5   & Bloque 3 --- Bienes públicos, bienes locales y federalismo fiscal & 2 & 20 \\
6--7   & Bloque 4 --- Elección pública & 2 & 20 \\
\midrule
\rowcolor{tonoA}\multicolumn{4}{l}{\cab{\textit{Parte II --- Gasto público, seguro social y redistribución}}} \\
8--10  & Bloque 5 --- Análisis del gasto público y evaluación social & 3 & 30 \\
11--12 & Bloque 6 --- Seguro social y pensiones & 2 & 20 \\
13--14 & Bloque 7 --- Salud, educación y redistribución & 2 & 20 \\
\midrule
\rowcolor{tonoA}\multicolumn{4}{l}{\cab{\textit{Parte III --- Tributación}}} \\
15--17 & Bloque 8 --- Incidencia y eficiencia de los impuestos & 3 & 30 \\
\rowcolor{tonoHielo}\textbf{18--20} & \textbf{Bloque 9 --- Tributación óptima} & \textbf{3} & \textbf{30} \\
21--23 & Bloque 10 --- Tributación aplicada & 3 & 30 \\
\midrule
\rowcolor{tonoA}\multicolumn{4}{l}{\cab{\textit{Parte IV --- Economía política y finanzas públicas en la práctica}}} \\
24--26 & Bloque 11 --- Economía política de la política fiscal & 3 & 30 \\
27--29 & Bloque 12 --- Política fiscal, deuda y sostenibilidad & 3 & 30 \\
30--31 & Bloque 13 --- Presupuesto, calidad del gasto y descentralización & 2 & 20 \\
32--33 & Bloque 14 --- Sistema tributario peruano y reforma fiscal & 2 & 20 \\
\midrule
\rowcolor{tonoHielo}\multicolumn{2}{l}{\textbf{Total}} & \textbf{33} & \textbf{330} \\
\bottomrule
\end{tabularx}
}

\medskip
\noindent\destacado{Ritmo supuesto.} Diez horas por semana, todas de contenido
nuevo: por ejemplo, cuatro sesiones de dos horas entre semana y una sesión de
dos horas el sábado para ejercicios y datos.

\medskip
\noindent\destacado{El reparto entre las partes.} Los fundamentos y el gasto
ocupan siete semanas cada uno; la tributación, nueve; la economía política y
las finanzas públicas en la práctica, diez. La tributación recibe más tiempo
porque reúne la teoría más exigente del plan (incidencia general y
tributación óptima) y la revisión de cada impuesto; la Parte IV, porque suma
el bloque de Drazen al estudio de las finanzas públicas peruanas.

\medskip
\noindent\destacado{El tramo largo.} El Bloque 9 (semanas 18--20) es el núcleo
analítico del plan: la tributación óptima exige la dualidad del consumidor, la
optimización con restricciones y, en el modelo de Mirrlees, la restricción de
compatibilidad de incentivos. Conviene repasar antes el Bloque 2 del sílabo
de microeconomía y el Bloque 7 del de economía matemática.

\medskip
\noindent\destacado{Hitos de control.} Conviene verificar el avance al terminar las
semanas 7, 14, 23 y 29. Si en alguno de esos cortes se acumula más de una semana de
retraso, se alarga el calendario y no se comprime el bloque siguiente.

\subsection*{Detalle semana a semana}

{\small\filasalternas
\renewcommand{\arraystretch}{1.1}\noindent\begin{tabularx}{\textwidth}{clXc}
\toprule
\rowcolor{tonoOscuro}\cab{Sem.} & \cab{Bloque} & \cab{Temas} & \cab{Horas} \\
\midrule
1  & 1 & 1.1 El sector público: alcance, tamaño y funciones; 1.2 Eficiencia, equidad y bienestar social & 10 \\
2  & 1 & 1.3 Herramientas empíricas: de la correlación a la causalidad; 1.4 Teorías del Estado: fallas de mercado y fallas de gobierno & 10 \\
\midrule
3  & 2 & 2.1 Teoría de las externalidades: Pigou, Coase y precios frente a cantidades; 2.2 Externalidades en la práctica: medio ambiente y tributación correctiva & 10 \\
\midrule
4  & 3 & 3.1 Bienes públicos: la regla de Samuelson y la provisión privada; 3.2 Mecanismos de provisión: Lindahl y Clarke--Groves & 10 \\
5  & 3 & 3.3 Bienes club, bienes públicos locales y la hipótesis de Tiebout; 3.4 Federalismo fiscal: asignación de funciones y transferencias & 10 \\
\midrule
6  & 4 & 4.1 Reglas de votación: unanimidad, mayoría y votante mediano; 4.2 Democracia representativa: competencia electoral y paradoja del voto & 10 \\
7  & 4 & 4.3 Búsqueda de rentas, grupos de interés y burocracia; 4.4 El tamaño del Estado: hechos y explicaciones & 10 \\
\midrule
8  & 5 & 5.1 El análisis del gasto público y el proceso de políticas en el Perú; 5.2 Análisis costo-beneficio I: el método y sus fundamentos & 10 \\
9  & 5 & 5.3 ACB II: valoración de impactos en los mercados; 5.4 ACB III: descuento, tasa social de descuento y riesgo & 10 \\
10 & 5 & 5.5 ACB IV: valoración sin mercado y ponderaciones distributivas; 5.6 Análisis de impacto regulatorio y costo-efectividad & 10 \\
\midrule
11 & 6 & 6.1 Teoría del seguro social: información asimétrica y seguro óptimo; 6.2 Pensiones: reparto, capitalización y reforma & 10 \\
12 & 6 & 6.3 Pensiones en América Latina y el Perú; 6.4 Seguro de desempleo, discapacidad y flexiguridad & 10 \\
\midrule
13 & 7 & 7.1 Salud: fallas de mercado, seguros y sistemas de salud; 7.2 Educación: provisión pública, competencia y retornos & 10 \\
14 & 7 & 7.3 Distribución del ingreso, pobreza y programas de transferencias; 7.4 Economía pública conductual & 10 \\
\bottomrule
\end{tabularx}
}

\newpage
{\small\filasalternas
\renewcommand{\arraystretch}{1.1}\noindent\begin{tabularx}{\textwidth}{clXc}
\toprule
\rowcolor{tonoOscuro}\cab{Sem.} & \cab{Bloque} & \cab{Temas} & \cab{Horas} \\
\midrule
15 & 8 & 8.1 Sistemas tributarios: tipos, principios y estructura; 8.2 Incidencia en equilibrio parcial & 10 \\
16 & 8 & 8.3 Incidencia en equilibrio general: el modelo de Harberger; 8.4 Incidencia fiscal neta & 10 \\
17 & 8 & 8.5 Exceso de gravamen y costo marginal de los fondos públicos; 8.6 Impuestos y oferta de trabajo & 10 \\
\midrule
18 & 9 & 9.1 Impuestos indirectos óptimos: Ramsey y la eficiencia productiva; 9.2 Muchos consumidores y reforma tributaria & 10 \\
19 & 9 & 9.3 El impuesto óptimo a la renta: el modelo de Mirrlees; 9.4 Fórmulas con elasticidades y progresividad óptima & 10 \\
20 & 9 & 9.5 Tributación mixta: el teorema de Atkinson--Stiglitz y las críticas; 9.6 La tributación del capital & 10 \\
\midrule
21 & 10 & 10.1 El impuesto a la renta personal; 10.2 Impuestos al ahorro, a la riqueza y al riesgo & 10 \\
22 & 10 & 10.3 El impuesto a la renta de las empresas y la inversión; 10.4 Tributación internacional y competencia fiscal & 10 \\
23 & 10 & 10.5 El IVA y la imposición al consumo; 10.6 Evasión y administración tributaria & 10 \\
\midrule
24 & 11 & 11.1 Inconsistencia temporal y compromiso en la política fiscal; 11.2 Elecciones y ciclos políticos oportunistas & 10 \\
25 & 11 & 11.3 Ciclos partidarios, competencia y deuda estratégica; 11.4 Economía política de la redistribución & 10 \\
26 & 11 & 11.5 Inacción, demora y crisis; 11.6 Déficits, deuda e instituciones presupuestarias & 10 \\
\midrule
27 & 12 & 12.1 Las cuentas fiscales: presupuesto, déficit y deuda; 12.2 Sostenibilidad de la deuda y riesgos fiscales & 10 \\
28 & 12 & 12.3 El balance estructural; 12.4 Reglas e instituciones fiscales & 10 \\
29 & 12 & 12.5 Política fiscal anticíclica, multiplicadores y discrecionalidad; 12.6 La historia fiscal del Perú y el Marco Macroeconómico Multianual & 10 \\
\midrule
30 & 13 & 13.1 El sistema presupuestario y la gestión de las finanzas públicas; 13.2 Calidad del gasto, inversión pública y servicio civil & 10 \\
31 & 13 & 13.3 Empresas públicas, banca pública e inclusión financiera; 13.4 La descentralización fiscal en el Perú & 10 \\
\midrule
32 & 14 & 14.1 El sistema tributario peruano: estructura, rendimiento y gastos tributarios; 14.2 Recursos naturales y régimen fiscal minero & 10 \\
33 & 14 & 14.3 ¿Por qué recaudan poco los países en desarrollo?; 14.4 Reforma fiscal para el crecimiento y la equidad & 10 \\
\bottomrule
\end{tabularx}
}

%==============================================================================
\parte{Parte I --- Fundamentos}

%------------------------------------------------------------------------------
\bloque{Bloque 1 --- El sector público y las herramientas de la economía pública --- 2 semanas}
{Qué hace el Estado y con qué herramientas se lo estudia: el tamaño y las
funciones del sector público, los criterios de eficiencia y equidad, los
métodos empíricos y las razones ---y los límites--- de la intervención.
Cubre la introducción y la línea de base de 1ECO18, la sesión 1.1 de 1ECO81,
la semana 1 de 1ECO36 (2025-1) y la clase 1 de MIT 14.41.}

\tema{1.1 El sector público: alcance, tamaño y funciones}
\lecturas{\GR{1} \quad \SR{1} \quad \BG{1}}
\apoyo{\HM{2}}
\consultalib{\SR{2} \quad \AS{1} \quad \UR{I.1.3}}
\nota{Las cuatro preguntas de la economía pública según Gruber (cuándo
intervenir, cómo, con qué efectos y por qué el Estado hace lo que hace), las
funciones de asignación, distribución y estabilización, y el sector público
peruano en cifras: ingresos, gastos, deuda y niveles de gobierno.}

\tema{1.2 Eficiencia, equidad y bienestar social}
\lecturas{\GR{2} \quad \SR{3} \quad \BG{3}}
\apoyo{\HM{13}}
\consultalib{\SR{7} \quad \HM{6--7} \quad \MU{23} \quad \MU{25} \quad Feldman y Serrano (2006), caps.~9--10}
\nota{Maximización restringida, equilibrio y bienestar; los teoremas del
bienestar y la imposibilidad de los impuestos de suma fija personalizados;
utilitarismo, Rawls y comparabilidad interpersonal. Las funciones de
bienestar y el teorema de Arrow ya se vieron en los temas 11.4 y 11.5 del
sílabo de microeconomía; aquí se repasan con la vista puesta en la
tributación óptima (Bloque 9).}

\tema{1.3 Herramientas empíricas: de la correlación a la causalidad}
\lecturas{\GR{3} \quad \BG{10}}
\articulos{Gertler et al.\ (2016), caps.~1--3}
\consultalib{Vásquez (ed.) (2013), caps.~9--10}
\nota{Ensayos aleatorios, cuasi experimentos, diferencias en diferencias,
variables instrumentales y regresión discontinua aplicados a la política
pública. Los métodos están completos en el sílabo de microeconometría; aquí
basta con saber leer críticamente la evidencia que citan los textos.}

\tema{1.4 Teorías del Estado: fallas de mercado y fallas de gobierno}
\lecturas{\SR{4} \quad \BG{2} \quad \HM{3}}
\articulos{Banco Mundial (1997), caps.~1--2}
\consultalib{\AS{10--11} \quad \UR{I.2}}
\consulta{Stiglitz (1998) \quad Tanzi (2000) \quad Buchanan y Musgrave (1999) \quad Mazzucato (2021)}
\nota{Por qué fallan los mercados (competencia imperfecta, bienes públicos,
externalidades, mercados incompletos e información imperfecta) y por qué
también falla el Estado; las teorías del crecimiento del sector público
(Wagner, Baumol) y del gobierno excesivo (Niskanen, el Leviatán), que se
retoman en el tema 4.4. El Banco Mundial (1997) plantea el ajuste entre el papel del Estado y su
capacidad.}

%------------------------------------------------------------------------------
\bloque{Bloque 2 --- Externalidades --- 1 semana}
{Las externalidades y sus remedios privados y públicos, con su aplicación al
medio ambiente y a la tributación correctiva. Cubre el tema de externalidades
de 1ECO18, la sesión 1.4 de 1ECO81 y las clases 2 a 4 de MIT 14.41; amplía el
tema 11.1 del sílabo de microeconomía.}

\tema{2.1 Teoría de las externalidades: Pigou, Coase y precios frente a cantidades}
\lecturas{\GR{5} \quad \BG{4} \quad \HM{10}}
\consultalib{\SR{6} \quad \MU{2.6--2.12} \quad \UR{I.2.3} \quad Feldman y Serrano (2006), cap.~7}
\consulta{Baumol y Oates (1988), cap.~4 \quad Ostrom (2011)}
\nota{La ineficiencia del equilibrio con externalidades, el impuesto
pigouviano, los permisos negociables, la internalización y el teorema de
Coase; la elección entre precios y cantidades cuando hay incertidumbre
(Weitzman); los bienes comunes y su gobierno.}

\tema{2.2 Externalidades en la práctica: medio ambiente y tributación correctiva}
\lecturas{\GR{6.1--6.2} \quad \RB{11}}
\apoyo{\SA{10}}
\consultalib{Field y Field (2017), caps.~10--13 \quad \AS{14.5}}
\nota{La lluvia ácida y el comercio de emisiones, el calentamiento global y
los impuestos al carbono; los impuestos sobre los «males» en América Latina
(combustibles, tabaco, alcohol, emisiones) y el debate sobre el doble
dividendo. En el Perú cumple ese papel el impuesto selectivo al consumo.}

%------------------------------------------------------------------------------
\bloque{Bloque 3 --- Bienes públicos, bienes locales y federalismo fiscal --- 2 semanas}
{La provisión de bienes públicos y el reparto de funciones entre niveles de
gobierno: la regla de Samuelson, la provisión privada, los mecanismos de
revelación, los bienes club, el modelo de Tiebout y el federalismo fiscal.
Cubre el tema de bienes públicos de 1ECO18, las sesiones 1.2 y 1.3 y el
Módulo 2 (sesiones 2.1 a 2.3) de 1ECO81 y las clases 5 y 8 de MIT 14.41.}

\tema{3.1 Bienes públicos: la regla de Samuelson y la provisión privada}
\lecturas{\GR{7} \quad \BG{5.1--5.3} \quad \HM{8.1--8.4} \quad \HM{8.8--8.9}}
\articulos{Samuelson (1955)}
\consultalib{\SR{5} \quad \MU{2.1--2.5} \quad \MU{2.13} \quad \DR{9} \quad \UR{I.2.2}}
\consulta{Bergstrom, Blume y Varian (1986) \quad Oakland (1987)}
\nota{Rivalidad y exclusión, la condición de Samuelson, el problema del
polizón y la provisión privada: equilibrio de Nash con contribuciones
voluntarias, neutralidad frente a las transferencias, campañas de
recaudación y la evidencia experimental.}

\tema{3.2 Mecanismos de provisión: Lindahl y Clarke--Groves}
\lecturas{\BG{5.4--5.5} \quad \HM{8.5--8.7}}
\apoyo{\AS{16.1--16.6}}
\consultalib{\MU{8.1} \quad \UR{II.1.1} \quad Feldman y Serrano (2006), cap.~8}
\consulta{Laffont (1989), caps.~1--2}
\nota{El equilibrio de Lindahl, por qué nadie revela su disposición a pagar y
cómo el mecanismo pivotal de Clarke--Groves logra que se revele. Atkinson y
Stiglitz añaden la distribución y los bienes privados provistos públicamente.
El diseño de mecanismos en general está en los temas 13.3 a 13.5 del sílabo
de microeconomía.}

\tema{3.3 Bienes club, bienes públicos locales y la hipótesis de Tiebout}
\lecturas{\HM{9} \quad \GR{10.2}}
\articulos{Tiebout (1956) \quad Buchanan (1965)}
\consultalib{\AS{17.1--17.4} \quad \MU{9.1--9.5}}
\consulta{Berglas y Pines (1981)}
\nota{El tamaño óptimo de un club y la exclusión, los bienes públicos locales
y el «voto con los pies»; las pruebas empíricas, en particular la
capitalización de impuestos y servicios locales en el precio de la vivienda.}

\tema{3.4 Federalismo fiscal: asignación de funciones y transferencias}
\lecturas{\GR{10.1} \quad \GR{10.3} \quad \HM{19} \quad \NC{8} (diap.~1--10)}
\apoyo{\MU{10.1} \quad \MU{10.3}}
\articulos{Hines y Thaler (1995)}
\consultalib{\SR{26--27} \quad \AS{17.5} \quad \HM{20.3--20.4} \quad Oates (1999)}
\consulta{Musgrave y Musgrave (1992) \quad Shah (2006)}
\nota{El teorema de la descentralización, el dilema entre eficiencia y
estabilización, la rendición de cuentas y el riesgo compartido; las
transferencias en bloque, condicionadas y de contrapartida y el efecto
«papel matamoscas»: el dinero se queda donde cae. El caso peruano está en el
tema 13.4.}

%------------------------------------------------------------------------------
\bloque{Bloque 4 --- Elección pública --- 2 semanas}
{Cómo deciden los votantes, los políticos, los burócratas y los grupos de
interés: las reglas de votación, la competencia electoral, la búsqueda de
rentas, la burocracia y el tamaño del Estado. Mueller es el texto base; los
ciclos políticos y los déficits se estudian con Drazen en el Bloque 11.
Cubre el tema de elección social de 1ECO18, la parte de elección pública de
la sesión 1.3 de 1ECO81, la clase 7 de MIT 14.41 y el programa de elección
pública ECO1028 de Dublin City University (carpeta de economía política).}

\tema{4.1 Reglas de votación: unanimidad, mayoría y votante mediano}
\lecturas{\GR{9.1--9.2} \quad \BG{6} \quad \MU{4} \quad \MU{5.1--5.4}}
\apoyo{\HM{4}}
\consultalib{\MU{5.5--5.13} \quad \MU{6--7} \quad \MU{24} \quad \SR{9} \quad \UR{III.1.3} \quad Feldman y Serrano (2006), caps.~12--13}
\nota{La unanimidad de Wicksell y Lindahl y sus costos de decisión, la mayoría
óptima de Buchanan y Tullock, los ciclos de Condorcet y el votante mediano con
preferencias unimodales; el intercambio de votos y el control de la agenda. El
teorema de Arrow se estudió en el sílabo de microeconomía y \MU{24} lo
repasa.}

\tema{4.2 Democracia representativa: competencia electoral y paradoja del voto}
\lecturas{\GR{9.3} \quad \MU{11} \quad \MU{12.1--12.3} \quad \MU{14.1--14.4}}
\apoyo{Feddersen (2004)}
\consultalib{\MU{13} \quad \MU{14.5--14.7} \quad \PT{3.2--3.4} \quad Holcombe (1989)}
\nota{El modelo de Downs y sus límites (el conjunto no cubierto, candidatos
con preferencias, entrada de partidos), las pruebas del votante mediano, el
voto probabilístico y por qué la gente vota aunque su voto casi nunca sea
decisivo: el votante racional, el expresivo y el ético.}

\tema{4.3 Búsqueda de rentas, grupos de interés y burocracia}
\lecturas{\GR{9.4} \quad \MU{15} \quad \MU{16.1--16.2} \quad \MU{16.6}}
\articulos{Bardhan et al.\ (2024)}
\consultalib{\HM{5} \quad \MU{17} \quad \MU{20} \quad \PT{7.3} \quad \UR{III.1.4--1.5}}
\consulta{Olson (1992), caps.~I--II}
\nota{El juego de Tullock y la disipación de rentas; las rentas que crean la
regulación, los aranceles y otras actividades del Estado, y cuánto cuestan;
el burócrata que maximiza su presupuesto (Niskanen) y el Estado como
Leviatán; el clientelismo: cuándo los votantes premian las transferencias
focalizadas frente a los bienes públicos.}

\tema{4.4 El tamaño del Estado: hechos y explicaciones}
\lecturas{\MU{21} \quad \DR{14.1--14.3}}
\apoyo{\MU{22.4--22.5}}
\nota{El crecimiento del Estado en el siglo XX y sus explicaciones: el Estado
como proveedor de bienes públicos y como redistribuidor (Meltzer--Richard), la
ley de Wagner, el efecto Baumol, los grupos de interés, la burocracia y la
ilusión fiscal; y la relación entre el tamaño del Estado, la productividad y
el crecimiento.}

%==============================================================================
\parte{Parte II --- Gasto público, seguro social y redistribución}

%------------------------------------------------------------------------------
\bloque{Bloque 5 --- Análisis del gasto público y evaluación social --- 3 semanas}
{Cómo decide el Estado qué gastar y cómo evalúa lo que gasta: el marco del
análisis del gasto, el análisis costo-beneficio completo, el análisis de
impacto regulatorio y el de costo-efectividad. Cubre la Parte II de 1ECO18
(incluido el método de su trabajo del semestre) y la clase 6 de MIT 14.41.}

\tema{5.1 El análisis del gasto público y el proceso de políticas en el Perú}
\lecturas{\SR{10} \quad \BG{7} \quad \UR{II.2}}
\apoyo{Morón y Sanborn (2007), caps.~3--4}
\consultalib{\SR{12}}
\consulta{Zúñiga y Ramos (2024), unidad~2, cap.~3}
\nota{Las preguntas del análisis del gasto (necesidad, falla que corrige,
alternativas, eficiencia, distribución y equidad), las transferencias en
efectivo frente a los subsidios en especie y quiénes deciden en el Perú: el
Ejecutivo, el MEF, el Congreso, los gobiernos subnacionales y la burocracia.}

\tema{5.2 Análisis costo-beneficio I: el método y sus fundamentos}
\lecturas{\GR{8} \quad \BG{8} \quad \BO{1--2}}
\consultalib{\SR{11} \quad \BO{3} \quad \UR{II.1.2--1.3}}
\consulta{Barrantes (2018)}
\nota{Los pasos de un análisis costo-beneficio, el criterio de Kaldor--Hicks,
la disposición a pagar, el costo de oportunidad y quiénes cuentan
(\emph{standing}); la experiencia peruana con el análisis costo-beneficio de
normas.}

\tema{5.3 ACB II: valoración de impactos en los mercados}
\lecturas{\BO{5--8}}
\consultalib{\BO{4}}
\nota{Los excedentes en los mercados de productos y de insumos, los impactos
en mercados secundarios (cuándo contarlos y cuándo no) y la predicción y
monetización de los impactos. La estimación directa de la demanda
(\BO{4}) supone la econometría de los otros sílabos.}

\tema{5.4 ACB III: descuento, tasa social de descuento y riesgo}
\lecturas{\BO{9--10} \quad \AS{15.4}}
\consultalib{\BO{11--12}}
\nota{Valor presente, inflación y horizontes largos; los enfoques de la tasa
social de descuento (costo de oportunidad del capital, tasa social de
preferencia temporal, precio sombra del capital, tasas decrecientes). En el
Perú la tasa social de descuento la fija el MEF para los proyectos de
Invierte.pe (tema 13.2).}

\tema{5.5 ACB IV: valoración sin mercado y ponderaciones distributivas}
\lecturas{\BO{13} \quad \BO{15} \quad \BO{19}}
\consultalib{\BO{16--17} \quad \BO{20}}
\nota{El valor de existencia, los métodos de preferencias reveladas (precios
hedónicos, costo de viaje, gastos defensivos) y el análisis con pesos
distributivos, que conecta la evaluación de proyectos con la función de
bienestar social del tema 1.2.}

\tema{5.6 Análisis de impacto regulatorio y costo-efectividad}
\lecturas{\BG{9} \quad \BO{18}}
\articulos{OCDE (2008) \quad Ortiz de Zevallos y Guerra-García (1998)}
\consultalib{\BO{14}}
\consulta{Osinergmin (2016)}
\nota{Los análisis de necesidad, efectividad y proporcionalidad de una norma;
el análisis costo-efectividad y costo-utilidad cuando los beneficios no se
pueden monetizar. El método de Ortiz de Zevallos y Guerra-García es el que
usa 1ECO18 para evaluar proyectos de ley.}

%------------------------------------------------------------------------------
\bloque{Bloque 6 --- Seguro social y pensiones --- 2 semanas}
{Por qué el Estado asegura y cómo: la teoría del seguro social y su
aplicación a las pensiones, el desempleo y la discapacidad, con el sistema
previsional peruano. Cubre las clases 10 a 12 de MIT 14.41, la semana 12 de
1ECO36 y las clases 1 a 7 de 1ECO40.}

\tema{6.1 Teoría del seguro social: información asimétrica y seguro óptimo}
\lecturas{\GR{12} \quad \HM{12}}
\consultalib{Chetty y Finkelstein (2013)}
\nota{El valor del seguro con utilidad cóncava, la selección adversa y el
desmoronamiento de los mercados, el tamizado y la señalización, el riesgo
moral, el desplazamiento del autoseguro y la regla del seguro social óptimo
que equilibra el suavizamiento del consumo con el riesgo moral. Repasa el
Bloque 12 del sílabo de microeconomía.}

\tema{6.2 Pensiones: reparto, capitalización y reforma}
\lecturas{\GR{13} \quad \HM{22}}
\articulos{Barr y Diamond (2009)}
\consultalib{\HE{6.3--6.4}}
\consulta{Cipriani (2014) \quad Grech (2018)}
\nota{La seguridad social como seguro de longevidad y como transferencia
entre generaciones; sus efectos sobre el ahorro y la jubilación; la crisis de
los sistemas de reparto con el envejecimiento y los errores analíticos más
comunes en las reformas. El modelo de generaciones traslapadas con seguridad
social está en los temas 7.3 y 7.5 del sílabo de macroeconomía.}

\tema{6.3 Pensiones en América Latina y el Perú}
\lecturas{\NC{6}}
\articulos{BID (2019) \quad FUNDEF (2020) \quad Olivera (2020)}
\consulta{Arenas de Mesa (2019) \quad Bando, Galiani y Gertler (2020) \quad Bernal, Olivera y Suhrcke (2024) \quad Olivera e Iparraguirre (2022) \quad OCDE (2019)}
\nota{El sistema dual peruano (el Sistema Nacional de Pensiones, de reparto, y
el Sistema Privado, de capitalización), la baja cobertura por la
informalidad, las pensiones no contributivas (Pensión 65), las brechas de
género y el costo fiscal de las reformas. La Ley 32123 (2024), de
modernización del sistema previsional, es el marco vigente.}

\tema{6.4 Seguro de desempleo, discapacidad y flexiguridad}
\lecturas{\GR{14} \quad \SR{16}}
\apoyo{Kreiner y Svarer (2022)}
\consultalib{Cahuc, Carcillo y Zylberberg (2014), cap.~13}
\nota{El seguro de desempleo, el de discapacidad y el de accidentes de
trabajo: beneficios por suavizar el consumo frente a costos por riesgo moral;
el modelo danés, que combina despidos flexibles, prestaciones generosas y
obligaciones de búsqueda; las cuentas individuales de cesantía (la CTS
peruana) como alternativa.}

%------------------------------------------------------------------------------
\bloque{Bloque 7 --- Salud, educación y redistribución --- 2 semanas}
{Los grandes rubros del gasto social y la redistribución: salud, educación,
programas contra la pobreza y lo que la economía del comportamiento cambia en
el diseño de las intervenciones. Cubre las clases 9 y 13 a 16 de MIT 14.41,
los capítulos de gasto en la práctica de Stiglitz y Rosengard que lee 1ECO18
y la clase 10 de 1ECO40.}

\tema{7.1 Salud: fallas de mercado, seguros y sistemas de salud}
\lecturas{\GR{15} \quad \SR{13}}
\apoyo{McPake y Normand (2008), cap.~21}
\consultalib{\GR{16} \quad \HM{12.8} \quad McPake y Normand (2008), caps.~8, 19--20 y 22}
\nota{Por qué el mercado de la salud es peculiar (incertidumbre, información
asimétrica, demanda inducida), cuán generoso debe ser el seguro con los
pacientes y con los proveedores, y los modelos de sistemas: servicio nacional,
seguro social y seguro voluntario.}

\tema{7.2 Educación: provisión pública, competencia y retornos}
\lecturas{\GR{11} \quad \SR{14}}
\articulos{Dhaliwal et al.\ (2011)}
\consultalib{Cahuc, Carcillo y Zylberberg (2014), cap.~4}
\nota{Por qué el Estado financia y provee educación (externalidades,
restricciones de crédito, igualdad de oportunidades), los vales y la
competencia entre escuelas, los retornos privados y sociales, y la
comparación de intervenciones educativas por su costo-efectividad.}

\tema{7.3 Distribución del ingreso, pobreza y programas de transferencias}
\lecturas{\GR{17} \quad \HM{14} \quad \SA{8.1--8.4}}
\consultalib{\SR{15} \quad \SA{8.5--8.6} \quad Vásquez (ed.) (2013), caps.~1, 3 y 6}
\nota{La medición de la desigualdad (Gini, índice de Atkinson) y de la
pobreza (Foster--Greer--Thorbecke), el riesgo moral de los programas de
bienestar y cómo reducirlo (focalización, condicionalidades, requisitos), el
ingreso mínimo garantizado, el impuesto negativo, los subsidios a los
salarios bajos y el salario mínimo.}

\tema{7.4 Economía pública conductual}
\lecturas{\GR{6.3--6.4}}
\apoyo{Thaler y Sunstein (2003) \quad Madrian y Shea (2001)}
\articulos{Camerer et al.\ (2003) \quad Brown et al.\ (2008)}
\consultalib{Chetty, Looney y Kroft (2009) \quad DellaVigna (2009)}
\consulta{Allcott, Lockwood y Taubinsky (2018)}
\nota{El tabaquismo y las adicciones como «internalidades», el poder de las
opciones por defecto en el ahorro previsional, el enigma de la baja demanda de
rentas vitalicias y el paternalismo asimétrico o libertario: corregir sesgos
sin imponer costos a quien decide bien. Corresponde a los sesgos conductuales
de la sumilla de 1ECO18 y a las semanas 8 y 9 de ECO1028.}

%==============================================================================
\parte{Parte III --- Tributación}

%------------------------------------------------------------------------------
\bloque{Bloque 8 --- Incidencia y eficiencia de los impuestos --- 3 semanas}
{Quién paga realmente los impuestos y cuánto cuestan: los sistemas
tributarios, la incidencia en equilibrio parcial y general, la incidencia
fiscal neta, el exceso de gravamen y los efectos sobre la oferta de trabajo.
Cubre la Parte III de 1ECO18 salvo la tributación óptima, la sesión 3.1 de
1ECO81 y las clases 17 a 20 de MIT 14.41.}

\tema{8.1 Sistemas tributarios: tipos, principios y estructura}
\lecturas{\GR{18} \quad \SR{17} \quad \BG{11.1--11.3}}
\consultalib{\UR{IV.1.1--1.2}}
\nota{Tipos de impuestos, tasas medias y marginales, progresividad, equidad
horizontal y vertical, los principios del beneficio y de la capacidad de pago,
la base de Haig--Simons y la unidad de tributación; los atributos de un buen
sistema tributario: eficiencia, sencillez administrativa, flexibilidad,
transparencia y equidad.}

\tema{8.2 Incidencia en equilibrio parcial}
\lecturas{\GR{19.1--19.2} \quad \SR{18} \quad \SA{1.1}}
\apoyo{\HM{11.5--11.6}}
\consultalib{\BG{11.5} \quad \AS{6.1}}
\consulta{Fullerton y Metcalf (2002)}
\nota{Las tres reglas de la incidencia (la incidencia legal no importa, paga
el lado inelástico y la división depende de las elasticidades), los impuestos
a la planilla, la incidencia con competencia imperfecta y la diferencia entre
impuestos específicos y \emph{ad valorem}.}

\tema{8.3 Incidencia en equilibrio general: el modelo de Harberger}
\lecturas{\GR{19.3} \quad \SA{1.2} \quad \AS{6.2--6.4}}
\articulos{Harberger (1962)}
\consultalib{\AS{6.5--6.6} \quad \AS{7} \quad \UR{IV.2.2}}
\consulta{Fullerton y Ta (2017)}
\nota{El modelo de dos sectores y dos factores, la incidencia del impuesto a
la renta de las empresas y la generalización a cualquier impuesto; Atkinson y
Stiglitz extienden el análisis a mercados imperfectos y a economías sin
equilibrio de mercado. Urrunaga et al.\ exponen el modelo de Harberger en
castellano.}

\tema{8.4 Incidencia fiscal neta}
\lecturas{\AS{9.1--9.3} \quad \GR{19.4}}
\articulos{Piggott y Whalley (1987) \quad Lustig (2017)}
\consultalib{\AS{9.4--9.5}}
\consulta{Lustig, Pessino y Scott (2014)}
\nota{La incidencia distributiva conjunta de los impuestos y el gasto y sus
problemas de interpretación; la metodología del proyecto Compromiso con la
Equidad (CEQ) aplicada a América Latina, con el Perú entre los casos.}

\tema{8.5 Exceso de gravamen y costo marginal de los fondos públicos}
\lecturas{\GR{20.1} \quad \SR{19} \quad \SA{2.2--2.3}}
\articulos{Slemrod y Yitzhaki (2001)}
\consultalib{\HM{15.2} \quad \MU{22.1} \quad \UR{IV.1.3}}
\consulta{Dahlby (2008) \quad Auerbach y Hines (2002)}
\nota{El triángulo de Harberger, por qué la pérdida crece con el cuadrado de
la tasa, su medición con la variación equivalente y la curva de Laffer; el
costo marginal de los fondos públicos como puente entre la tributación y la
evaluación del gasto del Bloque 5.}

\tema{8.6 Impuestos y oferta de trabajo}
\lecturas{\GR{21} \quad \SA{2.1.1}}
\apoyo{\AS{2}}
\consultalib{\HM{16.2--16.3} \quad \HE{5.3}}
\nota{Efectos ingreso y sustitución, el margen extensivo y el intensivo, la
evidencia (elasticidades bajas para los hombres y más altas en la
participación de las mujeres casadas) y las políticas que promueven el
empleo, como el crédito tributario por ingresos laborales.}

%------------------------------------------------------------------------------
\bloque{Bloque 9 --- Tributación óptima --- 3 semanas}
{El núcleo normativo de la economía pública: cómo recaudar al menor costo de
eficiencia y con la equidad deseada, de Ramsey y Diamond--Mirrlees a
Mirrlees, las fórmulas de Saez, el teorema de Atkinson--Stiglitz y la
imposición del capital. Cubre las sesiones 3.2 a 3.4 de 1ECO81 y el tema de
tributación óptima de 1ECO18.}

\tema{9.1 Impuestos indirectos óptimos: Ramsey y la eficiencia productiva}
\lecturas{\HM{15.3--15.5} \quad \SA{3}}
\apoyo{\AS{12.1--12.3}}
\consultalib{\GR{20.2} \quad \SR{20} \quad \HM{15.9}}
\consulta{Myles (1995), cap.~4 \quad Diamond y Mirrlees (1971) \quad Stern (1987)}
\nota{La regla de Ramsey (reducción proporcional de las demandas
compensadas), la regla de la inversa de la elasticidad y sus casos
especiales; el teorema de eficiencia productiva de Diamond y Mirrlees: no
gravar los bienes intermedios. Los precios del sector público
(Ramsey--Boiteux) siguen la misma lógica.}

\tema{9.2 Muchos consumidores y reforma tributaria}
\lecturas{\HM{15.6--15.8} \quad \AS{12.4--12.5}}
\articulos{Ahmad y Stern (1984)}
\consulta{Myles (1995), cap.~6}
\nota{La regla de Ramsey con consumidores heterogéneos (Diamond), la teoría
de la reforma tributaria marginal, el óptimo inverso y las mejoras de Pareto;
la aplicación de Ahmad y Stern a los impuestos indirectos de la India.}

\tema{9.3 El impuesto óptimo a la renta: el modelo de Mirrlees}
\lecturas{\GR{20.3} \quad \HM{16.4--16.5} \quad \HM{17.1--17.5} \quad \SA{4}}
\consultalib{\AS{13}}
\consulta{Mirrlees (1971) \quad Sheshinski (1972) \quad Myles (1995), cap.~5 \quad Kaplow (2008), caps.~4--5}
\nota{El problema de información (el Estado observa el ingreso, no la
habilidad), el principio de revelación y la restricción de compatibilidad de
incentivos, el impuesto lineal óptimo y los resultados de Mirrlees: tasa
marginal nula en los extremos y tasas óptimas más bien planas en las
simulaciones. Repasa el tema 12.5 del sílabo de microeconomía.}

\tema{9.4 Fórmulas con elasticidades y progresividad óptima}
\lecturas{\HM{16.6--16.8}}
\articulos{Saez (2001) \quad Diamond y Saez (2011)}
\consulta{Piketty y Saez (2013) \quad Heathcote, Storesletten y Violante (2014) \quad Mertens (2013)}
\nota{La tasa marginal óptima en función de la elasticidad del ingreso
gravable, la forma de la distribución y los pesos sociales; la tasa óptima
para los ingresos más altos y el caso a favor de un impuesto progresivo.}

\tema{9.5 Tributación mixta: el teorema de Atkinson--Stiglitz y las críticas}
\lecturas{\SA{5} \quad \SA{7} \quad \HM{17.6--17.8}}
\articulos{Mankiw, Weinzierl y Yagan (2009)}
\consultalib{\AS{14.1--14.3}}
\consulta{Atkinson y Stiglitz (1976) \quad Jacobs y Boadway (2014) \quad Kaplow (2008), cap.~6 \quad Sørensen (2006) \quad Weinzierl (2014)}
\nota{El impuesto negativo como impuesto lineal; por qué, con separabilidad
débil entre consumo y ocio, los impuestos indirectos sobran cuando hay un
impuesto óptimo a la renta; las críticas (equidad horizontal, reformas,
administración y política) y las lecciones de la teoría para la política.}

\tema{9.6 La tributación del capital}
\lecturas{\SA{6}}
\apoyo{\KC{1--2}}
\consultalib{\AS{8} \quad \AS{14.4} \quad \HE{5.4} \quad \KC{3--5} \quad \DR{11.2}}
\consulta{Chamley (1986) \quad Kaplow (2008), cap.~9 \quad Myles (1995), cap.~13}
\nota{El resultado de tasa cero sobre el capital en el largo plazo
(Chamley--Judd), sus supuestos y sus excepciones con generaciones
traslapadas; el problema de credibilidad, que se estudia en el tema 11.1; la nueva finanza pública
dinámica, que sustituye el enfoque de Ramsey por el de Mirrlees dinámico y
justifica gravar el ahorro. La imposición del capital en el modelo de Ramsey
está en el tema 4.8 del sílabo de macroeconomía.}

%------------------------------------------------------------------------------
\bloque{Bloque 10 --- Tributación aplicada --- 3 semanas}
{Los impuestos uno por uno ---renta personal, ahorro y riqueza, renta
empresarial, tributación internacional e IVA--- y la evasión y la
administración tributaria, con la evidencia latinoamericana del BID. Cubre el
Módulo 4 de 1ECO81, las semanas 4, 5 y 7 de 1ECO36 y las clases 21 a 24 de
MIT 14.41.}

\tema{10.1 El impuesto a la renta personal}
\lecturas{\SR{22} \quad \RB{7} \quad \BG{11.4}}
\nota{La base imponible, las deducciones, la unidad de tributación y la
progresividad efectiva; por qué en América Latina el impuesto a la renta
personal es «un cascarón vacío»: pocos contribuyentes, rentas del capital que
escapan y deducciones que erosionan la base.}

\tema{10.2 Impuestos al ahorro, a la riqueza y al riesgo}
\lecturas{\GR{22} \quad \GR{23}}
\apoyo{\AS{3.1--3.2}}
\consultalib{\AS{3.3--3.5} \quad \AS{4} \quad \SA{2.1.2--2.1.3} \quad \SR{21}}
\nota{El ahorro y los impuestos en el modelo de ciclo de vida, las cuentas de
ahorro previsional con ventajas tributarias, los impuestos y la toma de
riesgos, las ganancias de capital, los impuestos a las herencias y el impuesto
a la propiedad.}

\tema{10.3 El impuesto a la renta de las empresas y la inversión}
\lecturas{\GR{24.1--24.5} \quad \RB{8}}
\apoyo{\AS{5.2}}
\consultalib{\SR{23} \quad \AS{5}}
\consulta{Auerbach, Devereux y Simpson (2008) \quad Gravelle y Smetters (2001)}
\nota{Por qué se grava a las sociedades, la estructura del impuesto, su
incidencia, el costo de uso del capital y los efectos sobre la inversión y el
financiamiento (Modigliani--Miller con impuestos, el sesgo hacia la deuda,
los dividendos); la competencia por la inversión mediante incentivos en
América Latina. El costo de uso del capital está en el tema 2.2 del sílabo de
macroeconomía.}

\tema{10.4 Tributación internacional y competencia fiscal}
\lecturas{\GR{24.6} \quad \SR{24} \quad \HM{20.2} \quad \NC{3} (diap.~32--38)}
\articulos{Cobham y Janský (2018)}
\consulta{OCDE (2013)}
\nota{Residencia frente a fuente, el crédito por impuestos pagados en el
exterior, los precios de transferencia y la planificación tributaria
agresiva; la competencia tributaria entre jurisdicciones y la estimación de la
pérdida de recaudación por elusión corporativa en cada país.}

\tema{10.5 El IVA y la imposición al consumo}
\lecturas{\GR{25} \quad \RB{9} \quad \SA{9}}
\consultalib{\SR{25}}
\consulta{Auerbach (2006) \quad Barreix, Bès y Roca (2010) \quad Kotlikoff, Lagarda y Marín (2023)}
\nota{Las equivalencias entre impuestos, el impuesto al consumo frente al
impuesto a la renta, el impuesto plano y el IVA en la práctica: base, tasas,
exoneraciones y devoluciones. El IVA personalizado que devuelve el impuesto a
los hogares pobres se discute en el tema 14.4.}

\tema{10.6 Evasión y administración tributaria}
\lecturas{\HM{18} \quad \RB{6} \quad \BG{11.6} \quad \NC{3} (diap.~25--31)}
\articulos{Hutton (2017)}
\consultalib{MEF (2024), recuadro~5}
\consulta{Slemrod y Yitzhaki (2002) \quad Sandmo (2004) \quad Keen y Slemrod (2017) \quad Silvani y Baer (1997)}
\nota{El modelo de Allingham y Sandmo, las auditorías y las sanciones, la
moral tributaria y las interacciones sociales; cómo se mide la brecha del IVA
y qué hace eficaz a una administración tributaria. Las diapositivas de
1ECO36 muestran la evasión del IGV y del impuesto a la renta de tercera
categoría en el Perú.}

%==============================================================================
\parte{Parte IV --- Economía política y finanzas públicas en la práctica}

%------------------------------------------------------------------------------
\bloque{Bloque 11 --- Economía política de la política fiscal --- 3 semanas}
{Por qué la política fiscal observada se aparta de la óptima: la
inconsistencia temporal y los mecanismos de compromiso, los ciclos políticos,
la economía política de la redistribución, las demoras de las reformas y el
sesgo hacia el déficit, con sus remedios institucionales. Drazen es el texto
base. El bloque abre la Parte IV porque usa la tributación del capital
(tema 9.6) y la redistribución (Bloque 7) y prepara las reglas fiscales del
Bloque 12. Cubre la semana 4 de ECO1028 y amplía los temas 2.8 y 6.6--6.7 del
sílabo de macroeconomía.}

\tema{11.1 Inconsistencia temporal y compromiso en la política fiscal}
\lecturas{\DR{4.1--4.3} \quad \DR{4.6} \quad \DR{5.1--5.3} \quad \DR{5.5}}
\apoyo{\PT{12.1--12.2}}
\consultalib{\DR{6} \quad \PT{11}}
\nota{La tentación de gravar el capital ya acumulado, la inconsistencia
temporal entendida como conflicto de intereses, el dilema entre compromiso y
flexibilidad, y las leyes, constituciones, delegaciones y estructuras fiscales
que permiten comprometerse. La versión monetaria (Barro--Gordon, reputación e
independencia del banco central) está en los temas 9.4 y 9.5 del sílabo de
macroeconomía.}

\tema{11.2 Elecciones y ciclos políticos oportunistas}
\lecturas{\DR{7.1--7.3} \quad \MU{19.4}}
\apoyo{Shi y Svensson (2006)}
\consultalib{\MU{19.1--19.2} \quad \PT{16}}
\nota{Cómo las elecciones disciplinan y a la vez distorsionan la política: el
ciclo oportunista de Nordhaus, los ciclos con expectativas racionales e
información asimétrica sobre la competencia del gobernante, y los ciclos
presupuestarios y su evidencia, que los encuentra mayores en los países en
desarrollo y donde hay menos votantes informados.}

\tema{11.3 Ciclos partidarios, competencia y deuda estratégica}
\lecturas{\DR{7.4--7.6} \quad \DR{7.9}}
\consultalib{\DR{7.7--7.8} \quad \MU{19.3} \quad \MU{19.5--19.6}}
\nota{Los ciclos partidarios (Hibbs y Alesina), la competencia del gobernante
y el voto prospectivo, el valor de las promesas de campaña y cómo un gobierno
puede atar las manos de su sucesor mediante la deuda, el puente con el
tema 11.6.}

\tema{11.4 Economía política de la redistribución}
\lecturas{\DR{8.1--8.3} \quad \DR{8.6--8.7} \quad \PT{6.1}}
\consultalib{\DR{8.4--8.5} \quad \PT{6.2--6.4} \quad \MU{3}}
\nota{El modelo de Meltzer y Richard y por qué la redistribución observada no
siempre crece con la desigualdad, las transferencias diferenciadas hacia los
grupos con poder político, la redistribución entre generaciones (la política
de la seguridad social) y la redistribución cuando hay movilidad social.
Complementa el análisis normativo del Bloque 7.}

\tema{11.5 Inacción, demora y crisis}
\lecturas{\DR{10.1--10.8}}
\consultalib{\DR{13.2--13.5}}
\nota{Por qué se demoran reformas que mejorarían el bienestar: los intereses
creados, la incertidumbre sobre quién gana (Fernandez y Rodrik), la guerra de
desgaste por la carga del ajuste (Alesina y Drazen), los modelos de propiedad
común y las crisis como catalizadoras de reformas. Es la base de la economía
política de la reforma fiscal (tema 14.4).}

\tema{11.6 Déficits, deuda e instituciones presupuestarias}
\lecturas{\DR{14.4--14.5} \quad \PT{13.1--13.4} \quad \MU{19.7}}
\articulos{Banco Mundial (2017), panorama general}
\consultalib{\PT{13.5}}
\nota{Las explicaciones políticas de los déficits persistentes (ilusión
fiscal, fondo común, deuda estratégica, inestabilidad política, estabilizaciones
demoradas) y las instituciones presupuestarias que los contienen: reglas
numéricas, procedimientos jerárquicos o colegiados y transparencia; la
gobernanza como condición para que funcionen. Parte de los temas 2.7 y 2.8
del sílabo de macroeconomía; las reglas fiscales peruanas se estudian en el
tema 12.4.}

%------------------------------------------------------------------------------
\bloque{Bloque 12 --- Política fiscal, deuda y sostenibilidad --- 3 semanas}
{La lectura de las cuentas fiscales: déficit, deuda y sostenibilidad, el
balance estructural, las reglas e instituciones fiscales, la política fiscal
anticíclica y la historia fiscal peruana, con el Marco Macroeconómico
Multianual y la opinión del Consejo Fiscal. Cubre las semanas 1 y 2 de
1ECO36 y complementa los temas 2.4 a 2.8 y el Bloque 13 del sílabo de
macroeconomía.}

\tema{12.1 Las cuentas fiscales: presupuesto, déficit y deuda}
\lecturas{\GR{4} \quad \NC{1} (diap.~1--22)}
\apoyo{\SR{28}}
\consultalib{\HE{7.2}}
\consulta{Arias (2025): «El Estado como botín», «Nos estamos consumiendo los ahorros fiscales» y «La sostenibilidad del déficit fiscal»}
\nota{Cómo se mide la posición fiscal (resultado económico y primario,
gobierno general frente a sector público no financiero, deuda bruta y neta),
los problemas de medición y por qué importan los déficits. La restricción
presupuestaria intertemporal y la equivalencia ricardiana están en los temas
2.4 y 2.5 del sílabo de macroeconomía.}

\tema{12.2 Sostenibilidad de la deuda y riesgos fiscales}
\lecturas{\HE{7.3} \quad \HE{7.6} \quad \NC{1} (diap.~30--33) \quad MEF (2024), secc.~7}
\apoyo{OCDE (2023a), secc.~3.4 y 3.6}
\consultalib{\HE{7.4--7.5} \quad \DR{12.8}}
\nota{La aritmética de la deuda, el saldo primario que la estabiliza y las
pruebas de tensión; los pasivos contingentes explícitos (garantías,
asociaciones público-privadas, demandas judiciales) e implícitos (empresas
públicas, desastres naturales); el impago soberano. La sostenibilidad y el
señoreaje están en el tema 2.6 del sílabo de macroeconomía.}

\tema{12.3 El balance estructural}
\lecturas{Granda, Liendo y Rojas (2012) \quad BCRP (2008) \quad \NC{1} (diap.~23--29) \quad MEF (2024), secc.~6.1}
\nota{Cómo se aísla del resultado fiscal el efecto del ciclo económico y de
los precios de los minerales para medir el impulso fiscal discrecional: la
brecha del producto, las elasticidades de los ingresos y los precios de largo
plazo de los metales. Es la base de la orientación de la política fiscal que
evalúan el MEF y el Consejo Fiscal.}

\tema{12.4 Reglas e instituciones fiscales}
\lecturas{\NC{2} \quad \MA{2}}
\apoyo{Butrón y Céspedes (2020)}
\articulos{BID (2012), cap.~1 \quad Bonato (2019)}
\consultalib{MEF (2024), recuadro~4 y secc.~6.5}
\consulta{Arias (2025): «Reglas fiscales, MMM y rebaja de calificación»}
\nota{Tipos de reglas (de déficit, de deuda, de gasto y de balance
estructural), cláusulas de escape y consejos fiscales; la secuencia peruana
desde la Ley de Prudencia y Transparencia Fiscal (1999) hasta el Decreto
Legislativo 1276 (2016) y sus modificaciones, la última por el Decreto
Legislativo 1621 (2024). La economía política de las reglas está en el
tema 11.6.}

\tema{12.5 Política fiscal anticíclica, multiplicadores y discrecionalidad}
\lecturas{Taylor (2000) \quad Auerbach (2009) \quad DeLong y Summers (2012)}
\apoyo{Fatás y Mihov (2013)}
\consultalib{\HE{4.4--4.5}}
\nota{Estabilizadores automáticos frente a medidas discrecionales, el nuevo
activismo fiscal tras la crisis de 2008, la política fiscal en una economía
deprimida con la tasa de interés en su límite inferior y el costo, para el
crecimiento, de una política fiscal volátil. El multiplicador en el modelo
neokeynesiano y la cota inferior cero están en el Bloque 13 del sílabo de
macroeconomía.}

\tema{12.6 La historia fiscal del Perú y el Marco Macroeconómico Multianual}
\lecturas{\MA{1} \quad \MA{5--7}}
\apoyo{MEF (2024), secc.~1--3, 6.2--6.3 y 10}
\consultalib{\MA{3.2--3.3} \quad \MA{4.2--4.3} \quad MEF (2026) \quad FMI (2026) \quad CEPAL (2026)}
\nota{De la dominancia fiscal de los años ochenta a la consolidación: el
papel de los precios de los minerales frente al de las leyes fiscales, los
beneficios de la responsabilidad fiscal y el golpe de la COVID-19. La
opinión del Consejo Fiscal (sección 10 del MMM) se lee junto con el marco;
conviene repetir el ejercicio con la edición vigente, el MMM 2027--2030 de
agosto de 2026.}

%------------------------------------------------------------------------------
\bloque{Bloque 13 --- Presupuesto, calidad del gasto y descentralización --- 2 semanas}
{Cómo se programa, ejecuta y evalúa el gasto en el Perú: el sistema
presupuestario, la revisión de pares de la OCDE, la calidad del gasto, la
inversión pública y el servicio civil, las empresas y la banca públicas y la
descentralización fiscal. Cubre las semanas 10, 11, 13 y 14 de 1ECO36, el
anexo de las notas de 1ECO18 y la sesión 2.4 de 1ECO81.}

\tema{13.1 El sistema presupuestario y la gestión de las finanzas públicas}
\lecturas{Banco Mundial (2011) \quad \BG{anexo} \quad \NC{5} (diap.~1--22) \quad OCDE (2023b)}
\apoyo{OCDE (2023a), cap.~1 y secc.~2.5--2.6}
\consultalib{OCDE (2023a), secc.~2.1--2.4 y 2.8--2.9 y cap.~3}
\consulta{Robinson y Last (2009)}
\nota{Los principios y las fases del proceso presupuestario (programación,
formulación, aprobación, ejecución y evaluación), el presupuesto por
resultados y el participativo, el SIAF y las modificaciones presupuestarias.
La Ley 28411 que describe Banco Mundial (2011) fue sustituida en 2018 por el
Decreto Legislativo 1440, dentro del marco de administración financiera del
Decreto Legislativo 1436; las diapositivas de 1ECO36 usan el marco vigente.}

\tema{13.2 Calidad del gasto, inversión pública y servicio civil}
\lecturas{\NC{5} (diap.~23--31) \quad OCDE (2023a), cap.~5}
\articulos{Banco Mundial (2025), caps.~3--4}
\consultalib{OCDE (2023a), cap.~4 \quad Vásquez (ed.) (2013), cap.~14}
\consulta{Armijo y Espada (2014) \quad Dal Bó, Finan y Rossi (2013) \quad Arias (2025): «Construyamos un Estado meritocrático»}
\nota{Las causas de la baja ejecución, la evaluación de la calidad del gasto,
el paso del SNIP a Invierte.pe, la brecha entre la programación y la ejecución
de la inversión y el régimen salarial del sector público. La revisión del
Banco Mundial (2025) es el diagnóstico más reciente.}

\tema{13.3 Empresas públicas, banca pública e inclusión financiera}
\lecturas{\SR{8} \quad \NC{7}}
\apoyo{Banco de la Nación (2026)}
\articulos{Pérez Caldentey y Titelman (2018), pp.~37--60}
\consultalib{\AS{15.1--15.3} \quad \DR{13.7}}
\nota{Por qué produce el Estado (monopolios naturales, objetivos sociales),
la eficiencia de las empresas públicas frente a las privadas, la
privatización y la regulación; el Banco de la Nación como agente financiero
del Estado (pagaduría, recaudación, pago de programas sociales) y como
instrumento de inclusión financiera donde no llega la banca privada.}

\tema{13.4 La descentralización fiscal en el Perú}
\lecturas{\NC{8} (diap.~11--30) \quad \RB{5}}
\articulos{CEPAL (2017), pp.~93--116 \quad Mendoza (2023)}
\consultalib{OCDE (2023a), secc.~2.7}
\consulta{Bahl y Martínez-Vázquez (2006) \quad Banco Mundial (2010) \quad Banco Mundial (2021) \quad Arias (2025): «Repensando la descentralización»}
\nota{Las transferencias intergubernamentales (el canon, los recursos
ordinarios y el Foncomún), la baja recaudación propia de los gobiernos
regionales y locales (el impuesto predial), la pereza fiscal y el balance
del proceso iniciado en 2002.}

%------------------------------------------------------------------------------
\bloque{Bloque 14 --- Sistema tributario peruano y reforma fiscal --- 2 semanas}
{El cierre aplicado: la estructura y el rendimiento del sistema tributario
peruano, los gastos tributarios, el régimen fiscal minero, por qué los
países en desarrollo recaudan poco y qué reforma fiscal haría falta. Cubre las
semanas 3, 6 y 8 de 1ECO36, el tema de tributación en el Perú de 1ECO18 y sus
lecturas de control sobre tributación y desarrollo.}

\tema{14.1 El sistema tributario peruano: estructura, rendimiento y gastos tributarios}
\lecturas{\NC{3} (diap.~1--24) \quad \BG{12}}
\articulos{CEPAL (2017), cap.~1 \quad CEPAL (2019)}
\consultalib{MEF (2024), recuadro~6}
\consulta{MEF (2022) \quad Arias (2018) \quad Arias (2025): «¿A quién benefician los gastos tributarios?» y «Más privilegios tributarios»}
\nota{Las fuentes del presupuesto, el IGV, el ISC, los aranceles, el impuesto
a la renta y el predial; el nivel y la estructura de la carga tributaria
frente a América Latina y la OCDE; las brechas de política y de cumplimiento
y los gastos tributarios (exoneraciones y regímenes especiales) y su costo.}

\tema{14.2 Recursos naturales y régimen fiscal minero}
\lecturas{\RB{10} \quad \NC{3} (diap.~39--51) \quad \MA{3.1} \quad \MA{4.1}}
\apoyo{FMI (2012)}
\consultalib{Dammert (2021), secc.~1.3 y 15.2}
\consulta{Arias (2025): «La renta minera» y «La exitosa reforma del régimen fiscal minero»}
\nota{La renta económica de los recursos y el diseño de su tributación
(regalías, impuestos sobre utilidades, impuestos a la renta del recurso); la
reforma minera peruana de 2011 (regalía sobre la utilidad operativa,
impuesto especial y gravamen especial a la minería), la carga fiscal efectiva
y la dependencia de la recaudación respecto de los precios de los minerales.}

\tema{14.3 ¿Por qué recaudan poco los países en desarrollo?}
\lecturas{\RB{1} \quad \RB{4}}
\articulos{Besley y Persson (2014) \quad Tanzi y Zee (2001)}
\consultalib{\RB{2--3} \quad \MU{12.5}}
\consulta{Stiglitz (2009)}
\nota{La capacidad fiscal como inversión del Estado, la informalidad, la
estructura productiva y la legitimidad; los mitos sobre la tributación en
América Latina y las recomendaciones de política tributaria para países en
desarrollo.}

\tema{14.4 Reforma fiscal para el crecimiento y la equidad}
\lecturas{\NC{4} \quad \RB{13}}
\articulos{Arias (2021) \quad Banco Mundial (2025), cap.~2}
\consultalib{\RB{12} \quad \AS{18}}
\consulta{Arias (2025): «La madre de todas las reformas» y «El impuesto predial y la importancia del catastro fiscal»}
\nota{Las propuestas para el Perú: ampliar la base del impuesto a la renta
personal y gravar mejor las rentas del capital, racionalizar las
exoneraciones y los regímenes especiales, fortalecer el predial y el
catastro, el IVA personalizado (tema 10.5), la facturación electrónica, la
calidad del gasto y el retorno a la sostenibilidad; y la economía política de
armar un programa de reformas (tema 11.5).}

%------------------------------------------------------------------------------
\newpage
\section*{Anexo A. De los capítulos a los bloques}
\addcontentsline{toc}{section}{Anexo A. De los capítulos a los bloques}

Capítulos de cada libro que entran en el plan y temas donde se leen;
«(c)» indica que se citan solo como consulta. Lo que queda fuera se detalla
en el Anexo B.

\newcommand{\mapa}[2]{%
\subsection*{#1}
{\small\filasalternas
\noindent\begin{tabularx}{\textwidth}{>{\raggedright\arraybackslash}p{1.1cm}X>{\raggedright\arraybackslash}p{3.4cm}}
\toprule
\rowcolor{tonoOscuro}\cab{Cap.} & \cab{Título} & \cab{Temas} \\
\midrule
#2
\bottomrule
\end{tabularx}
}}

\newcommand{\mapadoc}[2]{%
\subsection*{#1}
{\small\filasalternas
\noindent\begin{tabularx}{\textwidth}{X>{\raggedright\arraybackslash}p{3.4cm}}
\toprule
\rowcolor{tonoOscuro}\cab{Documento} & \cab{Temas} \\
\midrule
#2
\bottomrule
\end{tabularx}
}}

\mapa{A.1 Gruber, \emph{Public Finance and Public Policy} (3.ª ed.)}{%
1 & Why Study Public Finance? & 1.1  \\
2 & Theoretical Tools of Public Finance & 1.2  \\
3 & Empirical Tools of Public Finance & 1.3  \\
4 & Budget Analysis and Deficit Financing & 12.1  \\
5 & Externalities: Problems and Solutions & 2.1  \\
6 & Externalities in Action: Environmental and Health Externalities (6.1--6.2; 6.3--6.4) & 2.2; 7.4  \\
7 & Public Goods & 3.1  \\
8 & Cost-Benefit Analysis & 5.2  \\
9 & Political Economy (9.1--9.2; 9.3; 9.4) & 4.1; 4.2; 4.3  \\
10 & State and Local Government Expenditures (10.1, 10.3; 10.2) & 3.4; 3.3  \\
11 & Education & 7.2  \\
12 & Social Insurance: The New Function of Government & 6.1  \\
13 & Social Security & 6.2  \\
14 & Unemployment Insurance, Disability Insurance, and Workers' Compensation & 6.4  \\
15 & Health Insurance I: Health Economics and Private Health Insurance & 7.1  \\
16 & Health Insurance II: Medicare, Medicaid, and Health Care Reform & 7.1 (c)  \\
17 & Income Distribution and Welfare Programs & 7.3  \\
18 & Taxation in the United States and Around the World & 8.1  \\
19 & The Equity Implications of Taxation: Tax Incidence (19.1--19.2; 19.3; 19.4) & 8.2; 8.3; 8.4  \\
20 & Tax Inefficiencies and Their Implications for Optimal Taxation (20.1; 20.2; 20.3) & 8.5; 9.1 (c); 9.3  \\
21 & Taxes on Labor Supply & 8.6  \\
22--23 & Taxes on Savings; Taxes on Risk Taking and Wealth & 10.2  \\
24 & Corporate Taxation (24.1--24.5; 24.6) & 10.3; 10.4  \\
25 & Fundamental Tax Reform & 10.5  \\
}

\mapa{A.2 Stiglitz y Rosengard, \emph{Economics of the Public Sector} (4.ª ed.)}{%
1 & Defining Public Sector Responsibilities & 1.1  \\
2 & Measuring Public Sector Size & 1.1 (c)  \\
3 & Market Efficiency & 1.2  \\
4 & Market Failure & 1.4  \\
5 & Public Goods and Publicly Provided Private Goods & 3.1 (c)  \\
6 & Externalities and the Environment & 2.1 (c)  \\
7 & Efficiency and Equity & 1.2 (c)  \\
8 & Public Production of Goods and Services & 13.3  \\
9 & Public Choice & 4.1 (c)  \\
10 & Framework for Analysis of Expenditure Policy & 5.1  \\
11 & Evaluating Public Expenditure & 5.2 (c)  \\
12 & Defense, Research, and Technology & 5.1 (c)  \\
13 & Health Care & 7.1  \\
14 & Education & 7.2  \\
15 & Welfare Programs and the Redistribution of Income & 7.3 (c)  \\
16 & Social Insurance & 6.4  \\
17 & Introduction to Taxation & 8.1  \\
18 & Tax Incidence & 8.2  \\
19 & Taxation and Economic Efficiency & 8.5  \\
20 & Optimal Taxation & 9.1 (c)  \\
21 & Taxation of Capital & 10.2 (c)  \\
22 & The Personal Income Tax & 10.1  \\
23 & The Corporation Income Tax & 10.3 (c)  \\
24 & A Student's Guide to Tax Avoidance & 10.4  \\
25 & Reform of the Tax System & 10.5 (c)  \\
26--27 & Intergovernmental Fiscal Relations; Subnational Taxes and Expenditures & 3.4 (c)  \\
28 & Fiscal Deficits and Government Debt & 12.1  \\
}

\mapa{A.3 Hindriks y Myles, \emph{Intermediate Public Economics} (borrador de 2004)}{%
2 & Government & 1.1  \\
3 & Theories of the Public Sector & 1.4  \\
4 & Voting & 4.1  \\
5 & Rent-Seeking & 4.3 (c)  \\
6--7 & Competitive Economies; Efficiency of Competition & 1.2 (c)  \\
8 & Public Goods (8.1--8.4, 8.8--8.9; 8.5--8.7) & 3.1; 3.2  \\
9 & Club Goods and Local Public Goods & 3.3  \\
10 & Externalities & 2.1  \\
11 & Imperfect Competition (11.5--11.6) & 8.2  \\
12 & Asymmetric Information (12.8) & 6.1; 7.1 (c)  \\
13 & Optimality and Comparability & 1.2  \\
14 & Inequality and Poverty & 7.3  \\
15 & Commodity Taxation (15.2; 15.3--15.5; 15.6--15.8; 15.9) & 8.5 (c); 9.1; 9.2; 9.1 (c)  \\
16 & Income Taxation (16.2--16.3; 16.4--16.5; 16.6--16.8) & 8.6 (c); 9.3; 9.4  \\
17 & Limits to Redistribution (17.1--17.5; 17.6--17.8) & 9.3; 9.5  \\
18 & Tax Evasion & 10.6  \\
19 & Fiscal Federalism & 3.4  \\
20 & Fiscal Competition (20.2; 20.3--20.4) & 10.4; 3.4 (c)  \\
22 & Social Security and Debt & 6.2  \\
}

\mapa{A.4 Atkinson y Stiglitz, \emph{Lectures on Public Economics}}{%
1 & Introduction: Public Economics & 1.1 (c)  \\
2 & Household Decisions, Income Taxation, and Labour Supply & 8.6  \\
3 & Taxation, Savings, and Decisions over Time (3.1--3.2) & 10.2  \\
4 & Taxation and Risk-Taking & 10.2 (c)  \\
5 & Taxation and the Firm (5.2) & 10.3  \\
6 & Tax Incidence: Simple Competitive Equilibrium Model (6.2--6.4) & 8.3  \\
7 & Tax Incidence: Departures from the Standard Model & 8.3 (c)  \\
8 & Taxation and Debt in a Growing Economy & 9.6 (c)  \\
9 & Distributional Effect of Taxation and Public Expenditure (9.1--9.3) & 8.4  \\
10--11 & Theories of the State and Public Economics; Normative Analysis & 1.4 (c)  \\
12 & The Structure of Indirect Taxation (12.1--12.3; 12.4--12.5) & 9.1; 9.2  \\
13 & The Structure of Income Taxation & 9.3 (c)  \\
14 & A More General Treatment of the Optimal Tax Problem & 9.5 (c); 9.6 (c); 2.2 (c)  \\
15 & Public Sector Pricing and Production (15.1--15.3; 15.4) & 13.3 (c); 5.4  \\
16 & Public Goods and Publicly Provided Private Goods & 3.2  \\
17 & Local Public Goods & 3.3 (c); 3.4 (c)  \\
18 & Public Economics: Theory and Policy & 14.4 (c)  \\
}

\mapa{A.5 Salanié, \emph{The Economics of Taxation}}{%
1 & Tax Incidence (1.1; 1.2) & 8.2; 8.3  \\
2 & Distortions and Welfare Losses (2.1.1; 2.2--2.3) & 8.6; 8.5  \\
3 & Indirect Taxation & 9.1  \\
4 & Direct Taxation & 9.3  \\
5 & Mixed Taxation & 9.5  \\
6 & The Taxation of Capital & 9.6  \\
7 & Criticisms of Optimal Taxation & 9.5  \\
8 & Low-Income Support (8.1--8.4) & 7.3  \\
9 & The Consumption Tax & 10.5  \\
10 & Environmental Taxation & 2.2  \\
}

\mapa{A.6 Mueller, \emph{Public Choice III}}{%
2 & The reason for collective choice -- allocative efficiency & 2.1 (c); 3.1 (c)  \\
3 & The reason for collective choice -- redistribution & 11.4 (c) \\
4 & The choice of voting rule & 4.1  \\
5 & Majority rule -- positive properties (5.1--5.4) & 4.1  \\
6--7 & Majority rule -- normative properties; Simple alternatives to majority rule & 4.1 (c)  \\
8 & Complicated alternatives to majority rule (8.1) & 3.2 (c)  \\
9 & Exit, voice, and disloyalty (9.1--9.5) & 3.3 (c)  \\
10 & Federalism (10.1, 10.3) & 3.4  \\
11 & Two-party competition -- deterministic voting & 4.2  \\
12 & Two-party competition -- probabilistic voting (12.1--12.3; 12.5) & 4.2; 14.3 (c)  \\
13 & Multiparty systems & 4.2 (c)  \\
14 & The paradox of voting (14.1--14.4) & 4.2  \\
15 & Rent seeking & 4.3  \\
16 & Bureaucracy (16.1--16.2, 16.6) & 4.3  \\
17 & Legislatures and bureaucracies & 4.3 (c)  \\
19 & Political competition and macroeconomic performance (19.4; 19.7; 19.1--19.3, 19.5--19.6) & 11.2; 11.6; 11.2--11.3 (c) \\
20 & Interest groups, campaign contributions, and lobbying & 4.3 (c)  \\
21 & The size of government & 4.4 \\
22 & Government size and economic performance (22.1; 22.4--22.5) & 8.5 (c); 4.4 \\
23--25 & Social welfare functions; The impossibility of a social ordering; A just social contract & 1.2 (c); 4.1 (c)  \\
}

\mapa{A.7 Drazen, \emph{Political Economy in Macroeconomics}}{%
4 & The Time-Consistency Problem (4.1--4.3, 4.6) & 11.1 \\
5 & Laws, Institutions, and Delegated Authority (5.1--5.3, 5.5) & 11.1 \\
6 & Credibility and Reputation & 11.1 (c) \\
7 & Elections and Changes of Policymakers (7.1--7.3; 7.4--7.6, 7.9; 7.7--7.8) & 11.2; 11.3; 11.3 (c) \\
8 & Redistribution (8.1--8.3, 8.6--8.7; 8.4--8.5) & 11.4; 11.4 (c) \\
9 & Public Goods & 3.1 (c) \\
10 & Inaction, Delay, and Crisis & 11.5 \\
11 & Factor Accumulation and Growth (11.2) & 9.6 (c) \\
12 & The International Economy (12.8) & 12.2 (c) \\
13 & Economic Reform and Transition (13.2--13.5; 13.7) & 11.5 (c); 13.3 (c) \\
14 & The Size of Government and the Number of Nations (14.1--14.3; 14.4--14.5) & 4.4; 11.6 \\
}

\mapa{A.8 Boardman et al., \emph{Cost--Benefit Analysis} (5.ª ed.)}{%
1--2 & Introduction to Cost--Benefit Analysis; Conceptual Foundations of Cost--Benefit Analysis & 5.2  \\
3--4 & Microeconomic Foundations of CBA; Direct Estimation of Demand Schedules & 5.2 (c); 5.3 (c)  \\
5--8 & Valuing Impacts in Output, Input and Secondary Markets; Predicting and Monetizing Impacts & 5.3  \\
9--10 & Discounting Future Impacts and Handling Inflation; The Social Discount Rate & 5.4  \\
11--12 & Dealing with Uncertainty; Risk, Option Price, and Option Value & 5.4 (c)  \\
13 & Existence Value & 5.5  \\
14 & Valuing Impacts from Observed Behavior: Experiments and Quasi-Experiments & 5.6 (c)  \\
15 & Valuing Impacts from Observed Behavior: Indirect Market Methods & 5.5  \\
16--17 & Contingent Valuation; Shadow Prices from Secondary Sources & 5.5 (c)  \\
18 & Cost--Effectiveness Analysis and Cost--Utility Analysis & 5.6  \\
19 & Distributionally Weighted CBA & 5.5  \\
20 & How Accurate Is CBA? & 5.5 (c)  \\
}

\mapa{A.9 Barrantes, Manrique y Glave, \emph{Economía pública} (2.ª ed.)}{%
1 & Una mirada al sector público del Perú en el siglo XXI & 1.1  \\
2 & El Estado y el sector público & 1.4  \\
3 & Fundamentos de la economía del bienestar & 1.2  \\
4 & Externalidades & 2.1  \\
5 & Bienes públicos (5.1--5.3; 5.4--5.5) & 3.1; 3.2  \\
6 & Elección social & 4.1  \\
7 & La acción del sector público en la economía & 5.1  \\
8 & Análisis costo-beneficio & 5.2  \\
9 & Análisis de impacto regulatorio & 5.6  \\
10 & Evaluación de impacto & 1.3  \\
11 & Teoría de la tributación (11.1--11.3; 11.4; 11.5; 11.6) & 8.1; 10.1; 8.2 (c); 10.6  \\
12 & Tributación en el Perú & 14.1  \\
Anexo & Presupuesto público en el Perú & 13.1  \\
}

\mapa{A.10 Persson y Tabellini, \emph{Political Economics}}{%
3 & Electoral Competition (3.2--3.4) & 4.2 (c)  \\
6 & General-Interest Politics (6.1; 6.2--6.4) & 11.4; 11.4 (c) \\
7 & Special-Interest Politics (7.3) & 4.3 (c)  \\
11--12 & Dynamic Policy Problems; Capital Taxation (12.1--12.2) & 11.1 (c); 11.1 \\
13 & Public Debt (13.1--13.4; 13.5) & 11.6; 11.6 (c) \\
16 & Electoral Cycles & 11.2 (c) \\
}

\mapa{A.11 Heer, \emph{Public Economics} y Kocherlakota, \emph{The New Dynamic Public Finance}}{%
HE 4 & Government Consumption (4.4--4.5) & 12.5 (c)  \\
HE 5 & Income Taxation (5.3; 5.4) & 8.6 (c); 9.6 (c)  \\
HE 6 & Pensions (6.3--6.4) & 6.2 (c)  \\
HE 7 & Public Debt (7.2; 7.3, 7.6; 7.4--7.5) & 12.1 (c); 12.2; 12.2 (c)  \\
KC 1--2 & Introduction; The Ramsey Approach and Its Problems & 9.6  \\
KC 3--5 & Dynamic Social Contracting; Dynamic Optimal Taxation; Intergenerational Taxation & 9.6 (c)  \\
}

\mapa{A.12 BID (2013), \emph{Recaudar no basta}}{%
1 & Los mitos al desnudo & 14.3  \\
2--3 & La política de la política tributaria; Impuestos para enfrentar los contratiempos & 14.3 (c)  \\
4 & Cuidado con la informalidad & 14.3  \\
5 & Más ingresos propios para más desarrollo local & 13.4  \\
6 & Aprovechando al máximo la administración tributaria & 10.6  \\
7 & El impuesto a la renta personal: un cascarón vacío & 10.1  \\
8 & El impuesto a la renta empresarial: el arte de competir por la inversión y aumentar la recaudación & 10.3  \\
9 & El impuesto al valor agregado: que sea lo que es & 10.5  \\
10 & Impuestos a las materias primas pensando en el futuro & 14.2  \\
11 & Gravar los «males» para proteger los bienes & 2.2  \\
12 & Los impuestos heterodoxos: lo bueno, lo malo y lo feo & 14.4 (c)  \\
13 & Cómo armar un programa de reformas tributarias & 14.4  \\
}

\mapa{A.13 Mendoza y Anastacio, \emph{La historia fiscal del Perú: 1980--2020}}{%
1 & Las finanzas públicas en la historia de América Latina y el Perú: ¿qué es lo que sabemos? & 12.6  \\
2 & El marco de análisis y las hipótesis de la investigación & 12.4  \\
3 & Hipótesis y evidencia: los principales hechos estilizados (3.1; 3.2--3.3) & 14.2; 12.6 (c)  \\
4 & Hipótesis y evidencia: principales resultados econométricos (4.1; 4.2--4.3) & 14.2; 12.6 (c)  \\
5--7 & Los beneficios de la responsabilidad fiscal; COVID-19; Conclusiones & 12.6  \\
}

\mapa{A.14 Urrunaga, Hiraoka y Risso, \emph{Fundamentos de economía pública}}{%
I & Intervención del Estado en una economía de mercado (1.3; 2; 2.2; 2.3) & 1.1 (c); 1.4 (c); 3.1 (c); 2.1 (c)  \\
II & Gasto público (1.1; 1.2--1.3; 2) & 3.2 (c); 5.2 (c); 5.1  \\
III & Elección pública (1.3; 1.4--1.5) & 4.1 (c); 4.3 (c)  \\
IV & Impuestos (1.1--1.2; 1.3; 2.2) & 8.1 (c); 8.5 (c); 8.3 (c)  \\
}

\mapa{A.15 Notas de clase de 1ECO36 (L.~A.~Arias)}{%
NC1 & La estabilidad macroeconómica: ingresos, gasto, déficit y deuda (diap.~1--22; 23--29; 30--33) & 12.1; 12.3; 12.2  \\
NC2 & Las reglas fiscales & 12.4  \\
NC3 & Teoría y práctica del sistema tributario peruano (diap.~1--24; 25--31; 32--38; 39--51) & 14.1; 10.6; 10.4; 14.2  \\
NC4 & Reforma fiscal para el crecimiento y la equidad & 14.4  \\
NC5 & El presupuesto público y la evaluación de la calidad del gasto (diap.~1--22; 23--31) & 13.1; 13.2  \\
NC6 & El sistema de pensiones & 6.3  \\
NC7 & La inclusión financiera y el rol del Banco de la Nación & 13.3  \\
NC8 & Descentralización fiscal (diap.~1--10; 11--30) & 3.4; 13.4  \\
}

\mapadoc{A.16 Informes y documentos de la carpeta}{%
OCDE (2023a), \emph{Gestión de las finanzas públicas en el Perú}: cap.~1, secc.~2.5--2.6; cap.~5; secc.~3.4 y 3.6 & 13.1; 13.2; 12.2  \\
OCDE (2023b), versión resumida de la revisión de pares & 13.1  \\
MEF (2024), \emph{Marco Macroeconómico Multianual 2025--2028}: secc.~7; 6.1; 1--3, 6.2--6.3 y 10 & 12.2; 12.3; 12.6  \\
Granda, Liendo y Rojas (2012), MEF; BCRP (2008), Notas de Estudios 51 & 12.3  \\
Butrón y Céspedes (2020), \emph{Moneda} 181 & 12.4  \\
Taylor (2000); Auerbach (2009); DeLong y Summers (2012); Fatás y Mihov (2013) & 12.5  \\
Banco Mundial (2011), \emph{El sistema presupuestario en el Perú} & 13.1  \\
Banco de la Nación (2026), «Banco de la Nación en 2025» & 13.3  \\
Morón y Sanborn (2007), \emph{Los desafíos del policymaking en el Perú} & 5.1  \\
Kreiner y Svarer (2022), \emph{Danish Flexicurity} & 6.4  \\
}

%------------------------------------------------------------------------------
\section*{Anexo B. Lo que no entra en el plan}
\addcontentsline{toc}{section}{Anexo B. Lo que no entra en el plan}

{\small\filasalternas
\noindent\begin{tabularx}{\textwidth}{>{\raggedright\arraybackslash}p{4.6cm}X}
\toprule
\rowcolor{tonoOscuro}\cab{Materia} & \cab{Capítulos y secciones} \\
\midrule
Métodos numéricos & Las secciones de simulación y calibración de Heer (modelos de Ramsey y de generaciones traslapadas resueltos numéricamente, \HE{4.4}, \HE{5.7}, \HE{6.5} y \HE{7.5}) y \KC{6}; corresponden a otro curso (Anexo D.3) \\
Macroeconomía ya cubierta & \HE{2--3} (modelos de Ramsey y de generaciones traslapadas), \HM{21} (eficiencia intertemporal) y \HM{23} (crecimiento): Bloques 4 y 7 del sílabo de macroeconomía \\
Organización industrial y regulación & \HM{11.1--11.4} y \HM{11.7--11.9} (competencia imperfecta y regulación de monopolios y oligopolios) y el programa 1ECO82; el tema 12.7 del sílabo de microeconomía cubre la regulación con información asimétrica \\
Instituciones políticas comparadas & \PT{4--5} y \PT{8--10} (agencia, políticos partidarios, reglas electorales y regímenes) y \MU{18} (dictadura): Bloque 6 del sílabo de macroeconomía \\
Política monetaria en Drazen & \DR{4.4--4.5} y \DR{5.4} (Barro--Gordon, señoreaje e independencia del banco central): temas 9.4 y 9.5 del sílabo de macroeconomía; \DR{6} solo como consulta \\
Otros capítulos de Drazen & \DR{2--3} (herramientas y mecanismos de decisión, ya cubiertos por los sílabos de micro y macro y por Mueller), \DR{11} salvo 11.2 (Bloque 6 del sílabo de macroeconomía), \DR{12} salvo 12.8 (economía internacional), \DR{13.6} y \DR{13.8} (transición de las economías planificadas) y \DR{14.6} (el número de naciones) (Anexo D.6) \\
Crecimiento y política monetaria en Persson y Tabellini & \PT{14} (crecimiento), \PT{15} y \PT{17} (credibilidad e instituciones de la política monetaria): Bloques 6 y 9 del sílabo de macroeconomía; \PT{16} (ciclos electorales), solo como consulta en 11.2 \\
Elección pública constitucional y normativa & \MU{26--29} (la constitución como contrato, derechos liberales y balance de la elección pública) (Anexo D.2) \\
Programas de los Estados Unidos & \GR{16} y las secciones institucionales de Gruber se leen como contexto; solo consulta \\
Partes de los programas PUCP & El cálculo actuarial, la demografía y la riqueza previsional de 1ECO40 (clases 3, 8 y 15) y las lecturas sobre género y pensiones (clase 12), salvo como consulta en 6.3 \\
\bottomrule
\end{tabularx}
}

\medskip
\noindent\destacado{Archivos de la carpeta que no se usaron.} De la carpeta de
economía pública: \emph{Teoría y política monetaria y bancaria} de Rivas
(\texttt{TMB.pdf}, duplicado en la carpeta de macroeconomía) y el documento
de política monetaria del Banco Central de Chile (\texttt{DPE\_38}), que
pertenecen a la teoría monetaria; los dos documentos de trabajo del BCRP
sobre modelos de proyección (2003 y 2022), la \emph{Guía metodológica de la
Nota Semanal} del BCRP, el manual de cuentas nacionales trimestrales del FMI,
las páginas guardadas de los cursos IMFx, el capítulo de programación
financiera y la nota del Banco Central de Chile sobre el stock de capital,
que corresponden a la macroeconomía aplicada; el programa de matemáticas del
BCRP (\texttt{MATE.pdf}) y la guía docente de métodos numéricos
(\texttt{2023\_33443}), ya considerados en el sílabo de economía matemática;
el programa de macroeconometría de la UdeSA; la nota sobre el impacto de la
COVID-19 en los pagos (sistema de pagos); el afiche del BCRP, que es una
imagen sin texto; los dos archivos comprimidos, las capturas de pantalla y
las imágenes guardadas.

\medskip
\noindent\destacado{Textos de 1ECO36 sin capa de texto.} Los doce capítulos de
Arias (2025), \emph{Horizontes fiscales y políticos}, de la carpeta de
1ECO36 son impresiones a PDF sin texto seleccionable: se citan solo como
consulta (PUCP), en los temas 12.1, 12.4, 13.2, 13.4, 14.1, 14.2 y 14.4. Los
otros dos capítulos que pide 1ECO36 (sobre protección social y la «pensión
por consumo») no están en la carpeta.

\medskip
\noindent\destacado{Versiones.} Gruber se usa en su 3.ª edición (2011) y
Hindriks y Myles en el borrador de 2004 (véase la nota de notación). El
archivo de \emph{Recaudar no basta} se llama «BID 2012», pero el libro es de
2013, como lo cita 1ECO36. La revisión de pares de la OCDE que 1ECO36 cita
como «OCDE (2022)» es el informe publicado en 2023 que está en la carpeta. El
libro de Mendoza y Anastacio está duplicado en la carpeta de macroeconomía;
se usa la copia de la carpeta de economía pública.

%------------------------------------------------------------------------------
\section*{Anexo C. Programas utilizados}
\addcontentsline{toc}{section}{Anexo C. Programas utilizados}

{\small\filasalternas
\noindent\begin{tabularx}{\textwidth}{>{\raggedright\arraybackslash}p{5.6cm}X}
\toprule
\rowcolor{tonoOscuro}\cab{Programa} & \cab{Temas del plan} \\
\midrule
1ECO18, Economía Pública (pregrado, 2025-1 y 2026-2, R.~Barrantes) & Introducción y línea de base, Bloque 1; Parte I, Bloques 2 a 4; Parte II, Bloque 5 y temas 7.1 a 7.3; Parte III, Bloque 8 y temas 9.1 y 9.3; tributación en el Perú, 14.1; lecturas de control, 1.4, 4.3, 11.6 y 14.3; sesgos conductuales, 7.4 \\
1ECO36, Finanzas Públicas (pregrado, 2025-1 y 2026-1, L.~A.~Arias) & Semanas 1 y 2, 12.1 y 12.4; semanas 3, 6 y 8, Bloque 14; semanas 4, 5 y 7, Bloque 10; semanas 10 y 11, 13.1 y 13.2; semana 12, 6.3; semana 13, 13.3; semana 14, 13.4 \\
1ECO81, Economía Pública (Maestría en Economía, 2025-1, E.~Beteta y F.~Vásquez) & Módulo 1, Bloques 1 a 3 y 4.1; módulo 2, 3.3, 3.4 y 13.4; módulo 3, Bloques 8 y 9; módulo 4, 10.3 y 10.6 \\
1ECO40, Pensiones y Seguridad Social (pregrado, 2024-2, J.~Olivera) & Clases 1 a 7, 6.2 y 6.3; clase 10, 7.4; clases 13 y 14, consulta en 6.3 \\
MIT OCW 14.41, Public Finance and Public Policy (otoño de 2024, J.~Gruber) & Clases 1 a 24: Bloques 1 a 8 y 10, en un orden parecido; conjuntos de problemas y exámenes en el Anexo D.1 \\
ECO1028, Public Choice (Dublin City University, 2026, B.~Gietner; carpeta de economía política) & Semanas 2, 3, 5 y 8, Bloque 4; semana 4, 11.2; semana 6, 3.4; semana 9, 7.4 \\
\bottomrule
\end{tabularx}
}

\medskip
\noindent\destacado{Programas de la carpeta que no se usaron}: 1ECO82 (Teoría de
la Regulación, 2026-1), que corresponde a la organización industrial, y
1ECO38 (Economía y Género, carpeta de bienestar). La semana 10 de ECO1028
(instituciones y desarrollo) ya está en el Bloque 6 del sílabo de
macroeconomía.

%------------------------------------------------------------------------------
\section*{Anexo D. Material opcional}
\addcontentsline{toc}{section}{Anexo D. Material opcional}

Todo lo de este anexo es opcional: no entra en el cómputo de horas ni modifica
el cronograma.

\subsection*{D.1 Ejercicios y exámenes de MIT 14.41}
El sitio de MIT OpenCourseWare de 14.41 (otoño de 2024) publica los videos de
las clases, los conjuntos de problemas y los exámenes, con sus soluciones.
Sirven de práctica para los Bloques 1 a 8 y 10; usan la 7.ª edición de
Gruber.

\subsection*{D.2 Elección pública constitucional y normativa}
\MU{26--29}: la constitución como contrato utilitarista, los derechos
liberales y el balance de lo que la elección pública ha aportado al estudio
de la política. Momento sugerido: tras el Bloque 4.

\subsection*{D.3 Finanzas públicas cuantitativas}
Las secciones de simulación de Heer (\HE{4.4}, \HE{5.7}, \HE{6.5} y \HE{7.5})
y \KC{6}: modelos calibrados de impuestos, pensiones y deuda. Requieren los
métodos numéricos que corresponden a otro curso.

\subsection*{D.4 Informes periódicos para seguir la coyuntura fiscal}
El Marco Macroeconómico Multianual (agosto) y el Informe de Actualización de
Proyecciones Macroeconómicas (abril) del MEF, con la opinión del Consejo
Fiscal; la sección de finanzas públicas del Reporte de Inflación del BCRP; la
Nota Tributaria de la SUNAT; el \emph{Fiscal Monitor} del FMI (abril y
octubre); el \emph{Panorama Fiscal de América Latina y el Caribe} de la
CEPAL; y las estadísticas tributarias de América Latina de la OCDE, la CEPAL,
el CIAT y el BID. Los capítulos de Arias (2025) de la carpeta comentan la
coyuntura fiscal peruana reciente.

\subsection*{D.5 Manuales de referencia}
Los cinco volúmenes del \emph{Handbook of Public Economics} (Auerbach y
Feldstein; Auerbach, Chetty, Feldstein y Saez), que citan 1ECO18 y 1ECO81;
Myles (1995), \emph{Public Economics}; Kaplow (2008), \emph{The Theory of
Taxation and Public Economics}; y Laffont (1989), \emph{Fundamentals of Public
Economics}.

\subsection*{D.6 Economía política de la macroeconomía abierta y de la transición}
\DR{12} (regímenes cambiarios, crisis, cooperación internacional, controles de
capital, deuda soberana y ayuda externa) y \DR{13} (reforma y transición de
las economías planificadas). Momento sugerido: tras el Bloque 11.

%------------------------------------------------------------------------------
\section*{Anexo E. Referencias}
\addcontentsline{toc}{section}{Anexo E. Referencias}

\subsection*{E.1 Textos del plan}
{\small
\begin{itemize}[leftmargin=1.4em,itemsep=1pt,label={}]
\item Atkinson, A. B. y J. E. Stiglitz (2015). \emph{Lectures on Public Economics}, edición actualizada. Princeton University Press.
\item Barrantes, R., S. Manrique y C. Glave (2023). \emph{Economía pública}, 2.ª ed. Material de Enseñanza 9, Departamento de Economía, PUCP.
\item Boardman, A. E., D. H. Greenberg, A. R. Vining y D. L. Weimer (2018). \emph{Cost--Benefit Analysis: Concepts and Practice}, 5.ª ed. Cambridge University Press.
\item Corbacho, A., V. Fretes Cibils y E. Lora (eds.) (2013). \emph{Recaudar no basta: los impuestos como instrumento de desarrollo}. Banco Interamericano de Desarrollo.
\item Drazen, A. (2000). \emph{Political Economy in Macroeconomics}. Princeton University Press (carpeta de economía política).
\item Gruber, J. (2011). \emph{Public Finance and Public Policy}, 3.ª ed. Worth Publishers.
\item Heer, B. (2019). \emph{Public Economics: The Macroeconomic Perspective}. Springer.
\item Hindriks, J. y G. D. Myles (2004). \emph{Intermediate Public Economics}. Borrador de junio de 2004; publicado por The MIT Press, 2006.
\item Kocherlakota, N. R. (2010). \emph{The New Dynamic Public Finance}. Princeton University Press.
\item Mendoza, W. y Y. Anastacio (2021). \emph{La historia fiscal del Perú: 1980--2020. Colapso, estabilización, consolidación y el golpe de la COVID-19}. Fondo Editorial PUCP.
\item Mueller, D. C. (2003). \emph{Public Choice III}. Cambridge University Press (carpeta de economía política).
\item Persson, T. y G. Tabellini (2000). \emph{Political Economics: Explaining Economic Policy}. The MIT Press (carpeta de economía política).
\item Salanié, B. (2003). \emph{The Economics of Taxation}. The MIT Press.
\item Stiglitz, J. E. y J. K. Rosengard (2015). \emph{Economics of the Public Sector}, 4.ª ed. W. W. Norton.
\item Urrunaga, R., T. Hiraoka y A. Risso (2014). \emph{Fundamentos de economía pública}. Apuntes de Estudio 43, Universidad del Pacífico.
\end{itemize}
}

\subsection*{E.2 Informes, documentos y diapositivas de la carpeta}
{\small
\begin{itemize}[leftmargin=1.4em,itemsep=1pt,label={}]
\item Arias, L. A. Notas de clase NC1 a NC8 (diapositivas). Finanzas Públicas, 1ECO36, PUCP.
\item Auerbach, A. J. (2009). Implementing the New Fiscal Policy Activism. \emph{American Economic Review}, 99(2), 543--549.
\item Banco de la Nación (2026). «Banco de la Nación en 2025», cap.~4 del informe institucional del banco.
\item Banco Mundial (2011). \emph{El sistema presupuestario en el Perú: notas para no especialistas} (N. Shack, «Sistema Nacional de Presupuesto Público»). Proyecto de Mejoramiento de los Servicios de Justicia, Lima.
\item BCRP (2008). Metodología de cálculo del resultado estructural. \emph{Notas de Estudios del BCRP}, 51.
\item Butrón, L. y N. Céspedes (2020). Dos décadas del Marco Macrofiscal peruano. \emph{Moneda}, 181, BCRP.
\item DeLong, J. B. y L. H. Summers (2012). Fiscal Policy in a Depressed Economy. \emph{Brookings Papers on Economic Activity}, primavera, 233--297.
\item Fatás, A. e I. Mihov (2013). Policy Volatility, Institutions, and Economic Growth. \emph{Review of Economics and Statistics}, 95(2), 362--376.
\item Granda, A., C. Liendo y C. Rojas (2012). \emph{El indicador de balance estructural: una herramienta analítica de la política fiscal}. Dirección General de Política Macroeconómica, MEF.
\item Kreiner, C. T. y M. Svarer (2022). Danish Flexicurity: Rights and Duties. CEBI Working Paper 16/22, Universidad de Copenhague.
\item MEF (2024). \emph{Marco Macroeconómico Multianual 2025--2028}. Ministerio de Economía y Finanzas, agosto de 2024.
\item Morón, E. y C. Sanborn (2007). \emph{Los desafíos del policymaking en el Perú: actores, instituciones y reglas de juego}. Documento de Trabajo 77, Centro de Investigación de la Universidad del Pacífico.
\item OCDE (2023a). \emph{Gestión de las finanzas públicas en el Perú}. OECD Publishing.
\item OCDE (2023b). \emph{Gestión de las finanzas públicas en el Perú: una revisión de pares de la OCDE. Mensajes principales} (versión resumida).
\item Taylor, J. B. (2000). Reassessing Discretionary Fiscal Policy. \emph{Journal of Economic Perspectives}, 14(3), 21--36.
\end{itemize}
}

\subsection*{E.3 Libros de consulta de la carpeta (sin horas)}
{\small
\begin{itemize}[leftmargin=1.4em,itemsep=1pt,label={}]
\item Arias, L. A. (2025). \emph{Horizontes fiscales y políticos: desafíos y retos para el nuevo gobernante del Perú}. Doce capítulos en la carpeta de 1ECO36, sin capa de texto.
\item Cahuc, P., S. Carcillo y A. Zylberberg (2014). \emph{Labor Economics}, 2.ª ed. The MIT Press (carpeta de economía del trabajo).
\item Dammert, A. (2021). \emph{Economía minera}. Fondo Editorial PUCP (carpeta de economía de los recursos naturales).
\item Feldman, A. M. y R. Serrano (2006). \emph{Welfare Economics and Social Choice Theory}, 2.ª ed. Springer (carpeta de economía del bienestar).
\item Field, B. C. y M. K. Field (2017). \emph{Environmental Economics: An Introduction}, 7.ª ed. McGraw-Hill (carpeta de economía de los recursos naturales).
\item McPake, B. y C. Normand (2008). \emph{Health Economics: An International Perspective}, 2.ª ed. Routledge (carpeta de economía de la salud).
\item Vásquez, E. (ed.) (2013). \emph{Inversión social: indicadores, bases de datos e iniciativas}. Universidad del Pacífico.
\end{itemize}
}

\subsection*{E.4 Referencias tomadas de los programas y lecturas externas}
{\small
\begin{itemize}[leftmargin=1.4em,itemsep=1pt,label={}]
\item Ahmad, E. y N. Stern (1984). The Theory of Reform and Indian Indirect Taxes. \emph{Journal of Public Economics}, 25(3), 259--298.
\item Allcott, H., B. Lockwood y D. Taubinsky (2018). Ramsey Strikes Back: Optimal Commodity Taxes and Redistribution in the Presence of Salience Effects. NBER Working Paper 24233.
\item Arenas de Mesa, A. (2019). \emph{Los sistemas de pensiones en la encrucijada: desafíos para la sostenibilidad en América Latina}. Libros de la CEPAL 159.
\item Arias, L. A. (2018). \emph{El Perú hacia la OCDE: la agenda pendiente para la política tributaria 2018--2021}. Grupo de Justicia Fiscal.
\item Arias, L. A. (2021). \emph{Política fiscal y tributaria frente a la pandemia global del coronavirus}.
\item Armijo, M. y M. V. Espada (2014). \emph{Calidad del gasto público y reformas institucionales en América Latina}. Serie Macroeconomía del Desarrollo 156, CEPAL.
\item Atkinson, A. B. y J. E. Stiglitz (1976). The Design of Tax Structure: Direct versus Indirect Taxation. \emph{Journal of Public Economics}, 6(1--2), 55--75.
\item Auerbach, A. J. (2006). The Choice between Income and Consumption Taxes: A Primer. NBER Working Paper 12307.
\item Auerbach, A. J., M. P. Devereux y H. Simpson (2008). Taxing Corporate Income. NBER Working Paper 14494.
\item Auerbach, A. J. y J. R. Hines (2002). Taxation and Economic Efficiency. En A. J. Auerbach y M. Feldstein (eds.), \emph{Handbook of Public Economics}, vol.~3, cap.~21. North-Holland.
\item Bahl, R. y J. Martínez-Vázquez (2006). Sequencing Fiscal Decentralization. World Bank Policy Research Working Paper 3914.
\item Banco Mundial (1997). \emph{Informe sobre el desarrollo mundial 1997: el Estado en un mundo en transformación}. Washington, D.\,C.
\item Banco Mundial (2010). \emph{El proceso de descentralización y su relación con la eficiencia del gasto público}. Informe 52885-PE.
\item Banco Mundial (2017). \emph{World Development Report 2017: Governance and the Law}. Washington, D.\,C.
\item Banco Mundial (2021). \emph{Repensar el futuro del Perú: notas de política para transformar al Estado en un gestor de bienestar y desarrollo}.
\item Banco Mundial (2025). \emph{Perú: revisión de las finanzas públicas}. Washington, D.\,C.
\item Bando, R., S. Galiani y P. Gertler (2020). The Effects of Noncontributory Pensions on Material and Subjective Well-Being. \emph{Economic Development and Cultural Change}, 68(4), 1233--1255.
\item Bardhan, P., S. Mitra, D. Mookherjee y A. Nath (2024). How Do Voters Respond to Welfare vis-à-vis Public Good Programs? Theory and Evidence of Political Clientelism. NBER Working Paper 32158.
\item Barr, N. y P. Diamond (2009). Reforming Pensions: Principles, Analytical Errors and Policy Directions. \emph{International Social Security Review}, 62(2), 5--29.
\item Barrantes, R. (2018). Derechos y rentabilidad del dinero público: ¿cuán útil es el análisis de beneficio-costo para garantizar el ejercicio de derechos a los ciudadanos? En R. Farfán, S. Mariani y C. O'Neill (eds.), \emph{Repensando las reglas de juego}, 57--84. Universidad del Pacífico.
\item Barreix, A., M. Bès y J. Roca (2010). \emph{El IVA personalizado: aumentando la recaudación y compensando a los más pobres}. BID.
\item Baumol, W. J. y W. E. Oates (1988). \emph{The Theory of Environmental Policy}, 2.ª ed. Cambridge University Press.
\item Berglas, E. y D. Pines (1981). Clubs, Local Public Goods and Transportation Models: A Synthesis. \emph{Journal of Public Economics}, 15(2), 141--162.
\item Bergstrom, T., L. Blume y H. Varian (1986). On the Private Provision of Public Goods. \emph{Journal of Public Economics}, 29(1), 25--49.
\item Bernal, N., J. Olivera y M. Suhrcke (2024). The Effects of Social Pensions on Nutrition-Related Health Outcomes of the Poor: Quasi-Experimental Evidence from Peru. \emph{Health Economics}, 33(5), 971--991.
\item Besley, T. y T. Persson (2014). Why Do Developing Countries Tax So Little? \emph{Journal of Economic Perspectives}, 28(4), 99--120.
\item BID (2012). \emph{Las instituciones fiscales del mañana}. Banco Interamericano de Desarrollo.
\item BID (2019). \emph{Diagnóstico del sistema de pensiones peruano y avenidas de reforma}. Nota Técnica IDB-TN-1776.
\item Bonato, L. (2019). Peru's Fiscal Framework: An Outsider Perspective.
\item Brown, J. R., J. R. Kling, S. Mullainathan y M. V. Wrobel (2008). Why Don't People Insure Late-Life Consumption? A Framing Explanation of the Under-Annuitization Puzzle. \emph{American Economic Review}, 98(2), 304--309.
\item Buchanan, J. M. (1965). An Economic Theory of Clubs. \emph{Economica}, 32(125), 1--14.
\item Buchanan, J. M. y R. A. Musgrave (1999). \emph{Public Finance and Public Choice: Two Contrasting Visions of the State}. The MIT Press.
\item Camerer, C., S. Issacharoff, G. Loewenstein, T. O'Donoghue y M. Rabin (2003). Regulation for Conservatives: Behavioral Economics and the Case for «Asymmetric Paternalism». \emph{University of Pennsylvania Law Review}, 151(3), 1211--1254.
\item CEPAL (2017). \emph{Consensos y conflictos en la política tributaria de América Latina}. CEPAL.
\item CEPAL (2019). Gastos tributarios y los objetivos de desarrollo sostenible. En \emph{Panorama Fiscal de América Latina y el Caribe 2019}, cap.~IV. CEPAL.
\item Chamley, C. (1986). Optimal Taxation of Capital Income in General Equilibrium with Infinite Lives. \emph{Econometrica}, 54(3), 607--622.
\item Chetty, R. y A. Finkelstein (2013). Social Insurance: Connecting Theory to Data. En A. J. Auerbach, R. Chetty, M. Feldstein y E. Saez (eds.), \emph{Handbook of Public Economics}, vol.~5, cap.~3. North-Holland.
\item Chetty, R., A. Looney y K. Kroft (2009). Salience and Taxation: Theory and Evidence. \emph{American Economic Review}, 99(4), 1145--1177.
\item Cipriani, G. P. (2014). Population Aging and PAYG Pensions in the OLG Model. \emph{Journal of Population Economics}, 27(1), 251--256.
\item Cobham, A. y P. Janský (2018). Global Distribution of Revenue Loss from Corporate Tax Avoidance: Re-estimation and Country Results. \emph{Journal of International Development}, 30(2), 206--232.
\item Dahlby, B. (2008). \emph{The Marginal Cost of Public Funds: Theory and Applications}. The MIT Press.
\item Dal Bó, E., F. Finan y M. A. Rossi (2013). Strengthening State Capabilities: The Role of Financial Incentives in the Call to Public Service. \emph{Quarterly Journal of Economics}, 128(3), 1169--1218.
\item DellaVigna, S. (2009). Psychology and Economics: Evidence from the Field. \emph{Journal of Economic Literature}, 47(2), 315--372.
\item Dhaliwal, I., E. Duflo, R. Glennerster y C. Tulloch (2011). Comparative Cost-Effectiveness Analysis to Inform Policy in Developing Countries: A General Framework with Applications for Education. J-PAL, MIT.
\item Diamond, P. A. y J. A. Mirrlees (1971). Optimal Taxation and Public Production I: Production Efficiency. \emph{American Economic Review}, 61(1), 8--27.
\item Diamond, P. y E. Saez (2011). The Case for a Progressive Tax: From Basic Research to Policy Recommendations. \emph{Journal of Economic Perspectives}, 25(4), 165--190.
\item Feddersen, T. J. (2004). Rational Choice Theory and the Paradox of Not Voting. \emph{Journal of Economic Perspectives}, 18(1), 99--112.
\item FMI (2012). \emph{Fiscal Regimes for Extractive Industries: Design and Implementation}. Fondo Monetario Internacional.
\item Fullerton, D. y G. E. Metcalf (2002). Tax Incidence. En A. J. Auerbach y M. Feldstein (eds.), \emph{Handbook of Public Economics}, vol.~4, cap.~26. North-Holland.
\item Fullerton, D. y C. L. Ta (2017). Public Finance in a Nutshell: A Cobb Douglas Teaching Tool for General Equilibrium Tax Incidence and Excess Burden. NBER Working Paper 23064.
\item FUNDEF (2020). \emph{Temas selectos de pensiones: tendencias, reformas recientes y otros aspectos relevantes}. Nota de Política Pública 2020-01, Fundación de Estudios Financieros.
\item Gertler, P. J., S. Martínez, P. Premand, L. B. Rawlings y C. M. J. Vermeersch (2016). \emph{La evaluación de impacto en la práctica}, 2.ª ed. BID y Banco Mundial.
\item Gravelle, J. y K. Smetters (2001). Who Bears the Burden of the Corporate Tax in the Open Economy? NBER Working Paper 8280.
\item Grech, A. G. (2018). What Makes Pension Reforms Sustainable? \emph{Sustainability}, 10(8).
\item Harberger, A. C. (1962). The Incidence of the Corporation Income Tax. \emph{Journal of Political Economy}, 70(3), 215--240.
\item Heathcote, J., K. Storesletten y G. L. Violante (2014). Optimal Tax Progressivity: An Analytical Framework. NBER Working Paper 19899.
\item Hines, J. R. y R. H. Thaler (1995). Anomalies: The Flypaper Effect. \emph{Journal of Economic Perspectives}, 9(4), 217--226.
\item Holcombe, R. G. (1989). The Median Voter Model in Public Choice Theory. \emph{Public Choice}, 61(2), 115--125.
\item Hutton, E. (2017). \emph{Modelo y metodología para la estimación de la brecha tributaria en relación con el impuesto al valor agregado}. Fondo Monetario Internacional.
\item Jacobs, B. y R. Boadway (2014). Optimal Linear Commodity Taxation under Optimal Non-linear Income Taxation. \emph{Journal of Public Economics}, 117, 201--210.
\item Kaplow, L. (2008). \emph{The Theory of Taxation and Public Economics}. Princeton University Press.
\item Keen, M. y J. Slemrod (2017). Optimal Tax Administration. IMF Working Paper 17/8.
\item Kotlikoff, L. J., G. Lagarda y G. Marín (2023). A Personalized VAT with Capital Transfers: A Reform to Protect Low-Income Households in Mexico. IDB Working Paper 1475.
\item Laffont, J.-J. (1989). \emph{Fundamentals of Public Economics}. The MIT Press.
\item Lustig, N. (2017). El impacto del sistema tributario y el gasto social en la distribución del ingreso y la pobreza en América Latina: una aplicación del marco metodológico del proyecto Compromiso con la Equidad (CEQ). \emph{El Trimestre Económico}, 84(335), 493--568.
\item Lustig, N., C. Pessino y J. Scott (2014). The Impact of Taxes and Social Spending on Inequality and Poverty in Argentina, Bolivia, Brazil, Mexico, Peru, and Uruguay: Introduction to the Special Issue. \emph{Public Finance Review}, 42(3), 287--303.
\item Madrian, B. C. y D. F. Shea (2001). The Power of Suggestion: Inertia in 401(k) Participation and Savings Behavior. \emph{Quarterly Journal of Economics}, 116(4), 1149--1187.
\item Mankiw, N. G., M. Weinzierl y D. Yagan (2009). Optimal Taxation in Theory and Practice. \emph{Journal of Economic Perspectives}, 23(4), 147--174.
\item Mazzucato, M. (2021). \emph{Mission Economy: A Moonshot Guide to Changing Capitalism}. Harper Business.
\item MEF (2022). \emph{Análisis del rendimiento de los tributos 2021}. Ministerio de Economía y Finanzas.
\item Mendoza, A. (2023). \emph{La descentralización fiscal 2017--2022}. Grupo Propuesta Ciudadana.
\item Mertens, K. (2013). Marginal Tax Rates and Income: New Time Series Evidence. NBER Working Paper 19171.
\item Mirrlees, J. A. (1971). An Exploration in the Theory of Optimum Income Taxation. \emph{Review of Economic Studies}, 38(2), 175--208.
\item Musgrave, R. A. y P. B. Musgrave (1992). \emph{Hacienda pública teórica y aplicada}, 5.ª ed. McGraw-Hill.
\item Myles, G. D. (1995). \emph{Public Economics}. Cambridge University Press.
\item Oakland, W. H. (1987). Theory of Public Goods. En A. J. Auerbach y M. Feldstein (eds.), \emph{Handbook of Public Economics}, vol.~2, cap.~9. North-Holland.
\item Oates, W. E. (1999). An Essay on Fiscal Federalism. \emph{Journal of Economic Literature}, 37(3), 1120--1149.
\item OCDE (2008). \emph{Introductory Handbook for Undertaking Regulatory Impact Analysis (RIA)}. OCDE.
\item OCDE (2013). \emph{Addressing Base Erosion and Profit Shifting}. OECD Publishing.
\item OCDE (2019). \emph{OECD Reviews of Pension Systems: Peru}. OECD Publishing.
\item Olivera, J. (2020). Reformando el sistema de pensiones en Perú: costo fiscal, nivel de pensiones, brecha de género y desigualdad. Documento de Trabajo 496, Departamento de Economía, PUCP.
\item Olivera, J. e Y. Iparraguirre (2022). The Gender Gap in Pension Savings: Evidence from Peru's Individual Capitalization System. LISER Working Paper 2022-06.
\item Olson, M. (1992). \emph{La lógica de la acción colectiva: bienes públicos y la teoría de grupos}. Limusa.
\item Ortiz de Zevallos, G. y G. Guerra-García (1998). \emph{Análisis costo-beneficio de las normas}. Instituto Apoyo y USAID.
\item Osinergmin (2016). \emph{Guía de Política Regulatoria N.º 1: guía metodológica para la realización de análisis de impacto regulatorio en Osinergmin}. Lima.
\item Ostrom, E. (2011). \emph{El gobierno de los bienes comunes: la evolución de las instituciones de acción colectiva}. Fondo de Cultura Económica.
\item Pérez Caldentey, E. y D. Titelman (eds.) (2018). \emph{La inclusión financiera para la inserción productiva y el papel de la banca de desarrollo}. CEPAL.
\item Piggott, J. y J. Whalley (1987). Interpreting Net Fiscal Incidence Calculations. \emph{Review of Economics and Statistics}, 69(4), 685--694.
\item Piketty, T. y E. Saez (2013). Optimal Labor Income Taxation. En A. J. Auerbach, R. Chetty, M. Feldstein y E. Saez (eds.), \emph{Handbook of Public Economics}, vol.~5, cap.~7. North-Holland.
\item Robinson, M. y D. Last (2009). \emph{A Basic Model of Performance-Based Budgeting}. Technical Notes and Manuals, Fondo Monetario Internacional.
\item Saez, E. (2001). Using Elasticities to Derive Optimal Income Tax Rates. \emph{Review of Economic Studies}, 68(1), 205--229.
\item Samuelson, P. A. (1955). Diagrammatic Exposition of a Theory of Public Expenditure. \emph{Review of Economics and Statistics}, 37(4), 350--356.
\item Sandmo, A. (2004). The Theory of Tax Evasion: A Retrospective View. Discussion Paper 31/04, Norwegian School of Economics.
\item Shah, A. (2006). \emph{Intergovernmental Fiscal Transfers: Principles and Practice}. Banco Mundial.
\item Sheshinski, E. (1972). The Optimal Linear Income-Tax. \emph{Review of Economic Studies}, 39(3), 297--302.
\item Shi, M. y J. Svensson (2006). Political Budget Cycles: Do They Differ across Countries and Why? \emph{Journal of Public Economics}, 90(8--9), 1367--1389.
\item Silvani, C. y K. Baer (1997). Designing a Tax Administration Reform Strategy: Experiences and Guidelines. IMF Working Paper 97/30.
\item Slemrod, J. y S. Yitzhaki (2001). Integrating Expenditure and Tax Decisions: The Marginal Cost of Funds and the Marginal Benefit of Projects. NBER Working Paper 8196.
\item Slemrod, J. y S. Yitzhaki (2002). Tax Avoidance, Evasion, and Administration. En A. J. Auerbach y M. Feldstein (eds.), \emph{Handbook of Public Economics}, vol.~3, cap.~22. North-Holland.
\item Sørensen, P. B. (2006). The Theory of Optimal Taxation: What Is the Policy Relevance? EPRU Working Paper.
\item Stern, N. (1987). The Theory of Optimal Commodity and Income Taxation: An Introduction. En D. Newbery y N. Stern (eds.), \emph{The Theory of Taxation for Developing Countries}. Oxford University Press.
\item Stiglitz, J. E. (1998). Distinguished Lecture on Economics in Government: The Private Uses of Public Interests: Incentives and Institutions. \emph{Journal of Economic Perspectives}, 12(2), 3--22.
\item Stiglitz, J. E. (2009). Development-Oriented Tax Policy. Initiative for Policy Dialogue Working Paper.
\item Tanzi, V. (2000). El papel del Estado y la calidad del sector público. \emph{Revista de la CEPAL}, 71, 7--22.
\item Tanzi, V. y H. Zee (2001). \emph{La política tributaria en los países en desarrollo}. Temas de Economía 27, Fondo Monetario Internacional.
\item Thaler, R. H. y C. R. Sunstein (2003). Libertarian Paternalism. \emph{American Economic Review}, 93(2), 175--179.
\item Tiebout, C. M. (1956). A Pure Theory of Local Expenditures. \emph{Journal of Political Economy}, 64(5), 416--424.
\item Weinzierl, M. (2014). The Promise of Positive Optimal Taxation: Normative Diversity and a Role for Equal Sacrifice. \emph{Journal of Public Economics}, 118, 128--142.
\item Zúñiga, H. y C. Ramos (2024). \emph{Manual de introducción al derecho}. Fondo Editorial de la Universidad del Pacífico.
\end{itemize}
}

\subsection*{E.5 Fuentes en línea consultadas}
{\small
\begin{itemize}[leftmargin=1.4em,itemsep=1pt,label={}]
\item CEPAL (2026). \emph{Panorama Fiscal de América Latina y el Caribe 2026} (LC/PUB.2026/4-P). Santiago.
\item FMI (2026). \emph{Fiscal Monitor}, abril de 2026. Washington, D.\,C.
\item MEF (2026). \emph{Marco Macroeconómico Multianual 2027--2030}. Ministerio de Economía y Finanzas, agosto de 2026.
\item MIT OpenCourseWare (2024). 14.41 \emph{Public Finance and Public Policy}, otoño de 2024 (J. Gruber).
\end{itemize}
}

\vspace{1.2cm}
\franjas[2pt]
\begin{center}\small\color{tonoPizarra}
Plan de lectura de elaboración propia. Los títulos de los capítulos
corresponden a las obras citadas y se reproducen únicamente con fines de
referencia bibliográfica. La asignación de capítulos y secciones se basa en
los índices de los libros de la carpeta; los temas de los cursos proceden de
los programas citados en el Anexo C.
\end{center}

\end{document}
